\documentclass[10pt,a4paper]{amsart}
\usepackage[margin=19mm]{geometry}
\usepackage{amsmath,amssymb,mathtools}
\usepackage[scr=rsfs,cal=euler]{mathalfa}
\usepackage{xfrac}
\usepackage[unicode,hidelinks,bookmarks=false]{hyperref}
\newcommand{\ie}{i.e.\ }
\newcommand{\cf}{cf.\ }
\newcommand{\resp}{respectively\ }
\newcommand{\Subsection}{\subsection}
\providecommand{\nobreakdash}{\nobreak\hbox{-}}
\setlength{\emergencystretch}{2em}
\multlinegap0pt
\usepackage{esint}
\usepackage[shortlabels]{enumitem}
\usepackage{amssymb,amsthm}
\setlength{\parindent}{0.25in}
%\setlength{\parskip}{0.2cm}
\setlength{\parskip}{0.5\baselineskip plus2pt minus2pt}
\numberwithin{equation}{section}
\newtheorem{theorem}{Theorem}[section]
\newtheorem{lemma}[theorem]{Lemma}
\newtheorem{proposition}[theorem]{Proposition}
\newtheorem{corollary}[theorem]{Corollary}
\newtheorem{claim}[theorem]{Claim}
\newtheorem{remark}[theorem]{Remark}
\newtheorem{definition}[theorem]{Definition}
\newtheorem{property}{Property}[section]
\newtheorem{example}[theorem]{Example}
\newtheorem{openproblem}[theorem]{Open Problem}
\def\E{{\mathbb E}}
\def\Z{{\mathbb Z}}
\def\R{{\mathbb R}}
\def\Q{{\mathbb Q}}
\def\N{{\mathbb N}}
\def\C{{\mathbb C}}

\def\cA{{\mathcal A}}
\def\cB{{\mathcal B}}
\def\cC{{\mathcal C}}
\def\cE{{\mathcal F}}
\def\cF{{\mathcal F}}
\def\cG{{\mathcal G}}
\def\cH{{\mathcal H}}
\def\cJ{{\mathcal J}}
\def\cK{{\mathcal K}}
\def\cL{{\mathcal L}}
\def\cM{{\mathcal M}}
\def\cO{{\mathcal O}}
\def\cQ{{\mathcal Q}}
\def\cR{{\mathcal R}}
\def\cT{{\mathcal T}}
\def\cU{{\mathcal U}}

\def\fM{{\mathfrak M}}

\def\a{\alpha}
\def\b{\beta}
\def\e{\varepsilon}
\def\d{\delta}
\def\D{\triangle}
\def\G{\Gamma}
\def\k{\kappa}
\def\l{\lambda}
\def\L{\Lambda}
\def\m{\mu}
\def\n{\nabla}
\def\p{\partial}
\def\r{\rho}
\def\s{\sigma}
\def\t{\tau}
\def\U{\Upsilon}
\def\w{\omega}
\def\W{\Omega}
\def\g{\gamma}
\def\z{\zeta}
\def\om{\omega}

\def\1{\left(}
\def\2{\right)}
\def\3{\left\{}
\def\4{\right\}}
\def\8{\infty}
\def\sm{\setminus}
\def\ss{\subseteq}
\def\cc{\subset\subset}

\def\uu{\overline{u}}
\def\ud{\underline{u}}
\def\vu{\overline{v}}
\def\vd{\underline{v}}

\newcommand{\mres}{\mathbin{\vrule height 1.6ex depth 0pt width
0.13ex\vrule height 0.13ex depth 0pt width 1.3ex}}

\DeclareMathOperator*{\dvg}{div}
\DeclareMathOperator*{\trace}{tr}
\DeclareMathOperator*{\diag}{diag}
\DeclareMathOperator*{\diam}{diam}
\DeclareMathOperator*{\supp}{supp}
\DeclareMathOperator*{\osc}{osc}
\DeclareMathOperator*{\dist}{dist}
\DeclareMathOperator*{\sign}{sign}
\DeclareMathOperator*{\length}{length}



\def\Lap{\Delta}
\def\grad{\nabla}

\def\En{\mathcal{E}}

\def\tautwo{\tau_2}
\def\tor{\mathrm{tor}}
\def\ei{\l_1}
\def\vpar{\eta}
\def\tpar{\mathfrak{T}}
\def\fv{\mathscr{V}_{\vpar}}
\def\fvbar{\mathscr{V}}
\def\vmax{v_{\text{max}}}
\def\vo{\mathscr{V}}

\def\err{\tau}
\def\Ep{\mathcal{F}_{\err}}
\def\minE{\En_{min}}
\def\Min{\mathcal{M}}
\def\compset{\mathcal{H}}
\def\Isom{G}
\def\Qb{Q}
\def\vol{{\textrm{vol}}}
\def\Om{{\Omega}}
\def\pa{{\partial}}
 \def\nl{\mathscr{h}_f}
\def\UP{\text{UP}}
\def\DO{\text{DO}}
\def\na{\nabla}
\def\baryEps{\bar{\e}}
\def\RHSclass{\mathscr{C}}
\def\TBF{q}
\def\ccdel{\bar{\delta}}
\DeclareFontFamily{U}{BOONDOX-scr}{\skewchar \font =45}
\DeclareFontShape{U}{BOONDOX-scr}{m}{n}{<-> s * [1.0] BOONDOX-r-cal}{}
\DeclareMathAlphabet{\mathscr}{U}{BOONDOX-scr}{m}{n}
\providecommand{\bysame}{\leavevmode\hbox to3em{\hrulefill}\thinspace}
\providecommand{\MR}{\relax\ifhmode\unskip\space\fi MR }
\providecommand{\MRhref}[2]{\href{http://www.ams.org/mathscinet-getitem?mr=#1}{#2}}
\providecommand{\href}[2]{#2}
\title[Quantitative Resolvent and Eigenfunction Stability for the Faber-Krahn Inequality]{Quantitative Resolvent and Eigenfunction Stability for the Faber-Krahn Inequality: stability of the first Dirichlet eigenfunction (Theorem 5.2)}
\author{Mark Allen, Dennis Kriventsov, and Robin Neumayer}
\date{Source: arXiv:2504.13053v2 (CC BY 4.0)}
\begin{document}
\setcounter{secnumdepth}{4}
\setcounter{tocdepth}{1}
\begin{abstract}
For a bounded open set $\Omega \subset \mathbb{R}^n$ with the same volume as the unit ball, the classical Faber-Krahn inequality says that the first Dirichlet eigenvalue $\lambda_1(\Omega)$ of the Laplacian is at least that of the unit ball $B$. We prove that the deficit $\lambda_1(\W)- \lambda_1(B)$ in the Faber-Krahn inequality controls the square of the distance between the resolvent operator $(-\Delta_\Omega)^{-1}$ for the Dirichlet Laplacian on $\Omega$ and the resolvent operator on the nearest unit ball $B(x_\Omega)$. The distance is measured by the operator norm from $L^{\infty}$ to $L^2$. As a main application, we show that the Faber-Krahn deficit  $\lambda_1(\Omega)- \lambda_1(B)$ controls the squared $L^2$ norm between $k$th eigenfunctions on $\W$ and $B(x_\Omega)$ for every $k \in \mathbb{N}.$ In both of these main theorems, the quadratic power is optimal.
\end{abstract}
\maketitle
\noindent\emph{Source and licence.} Quantitative Resolvent and Eigenfunction Stability for the Faber-Krahn Inequality. Version arXiv:2504.13053v2 (CC BY 4.0), distributed under the Creative Commons Attribution 4.0 International licence, http://creativecommons.org/licenses/by/4.0/, as recorded in the arXiv OAI licence metadata for this exact version and shown on the abstract page https://arxiv.org/abs/2504.13053v2. Full rights notice: licenses/arxiv-2504.13053-CC-BY-4.0.txt.\par\smallskip
\noindent\textit{Editorial scope.} Counted unit: the first-eigenfunction stability theorem, source label t:first (source0.tex lines 1006--1011), printed as Theorem 5.2 in the published version and as Theorem 5.2 here, delivered with its complete original author proof at source lines 1012--1071 as the single primary proof body. No proof marker is added and no mathematics is written, repaired or re-derived. The article is itself an amsart article (source0.tex line 2), so no publisher class is replaced or shipped: the wrapper supplies the pinned runtime preamble, and the article's own settings are copied verbatim from the source --- \textbackslash{}numberwithin\{equation\}\{section\} and the \textbackslash{}newtheorem declarations at source lines 26 and 29--38, the macro block at source lines 43--119 and 126--158, the \textbackslash{}parindent and \textbackslash{}parskip settings at source lines 15--17, and the article's own bibliography preamble commands. Consequently every printed number of the delivered unit is produced by the article's own counters and coincides with the published version: Sections 1--5 with their subsections, including the unnumbered subsections Discussion of the Proofs and Organization of the Paper, Theorems 1.1--1.3, Lemmas 2.1--2.6, Lemmas 3.1--3.4, Theorem 4.1 with Lemmas 4.2--4.4, Proposition 5.1, Theorem 5.2, and every numbered equation of the delivered sections. Required context retained uncounted: the whole body of the article from the Introduction (source line 173) to the line immediately preceding the counted statement (source line 1005), that is the Introduction with the three main theorems, the discussion of the proofs, the organization of the paper and the sharpness and literature discussion; the whole of the Preliminaries (torsional rigidity, qualitative stability, geometric measure theoretic preliminaries); the whole section on the truncated barycenter with Lemma 3.1 and the three simplifying lemmas; the whole section of resolvent estimates for nearly spherical sets; and the whole of Sections 5.1 and 5.2 up to the counted statement, including the complete proof of Theorem 1.2 (source lines 929--976) and the qualitative stability Proposition 5.1 (source lines 993--1002) on which the counted proof depends. Every definition, lemma, proposition, theorem, corollary, remark, numbered equation, footnote and citation of those source lines is retained, including the two footnotes at source lines 405 and 570. The article's abstract (source line 166) is reproduced verbatim in the wrapper front matter, and the article's own title and author lines are reproduced in the wrapper title block. Not part of this unit and not redistributed: Section 5.3 (the examples demonstrating sharpness), Sections 6 and 7 (free boundary estimates, parts 1 and 2) with the proof of Theorem 1.3, the closing section explaining the proof of Lemma 4.2 at source line 2585, the article's own bibliography listing and its title-page material. The delivered text cites 49 of the article's bibliography keys; those entries are transcribed verbatim from the e-print's own bibliography file v.7.bbl (classic \textbackslash{}bibitem format, 61 entries) in the article's own order with the article's printed bibliography numbers, so every printed citation number of the delivered unit is the published one, and the five entries that print their author as a dash in the published list (AKN1, AKN3, FMP09, KJ1, KJ3) keep that printed form because each one directly follows the entry it repeats, exactly as in the published list. Reference adaptation: four cross-reference keys, ten occurrences, point at material outside this unit and are delivered as \textbackslash{}href links to the abstract page of this exact version with the printed numbers retained (method reference\_only\_adaptation, recorded in delivery-evidence.json): \textbackslash{}ref\{ssec: examples\} prints 5.3 and occurs at source lines 214 and 254 (twice); \textbackslash{}ref\{lem: volume constraint lemma\} prints 7.15 and occurs at source line 264; \textbackslash{}ref\{sec: a priori and existence\} prints 6 and occurs at source lines 364 (twice) and 927; \textbackslash{}ref\{sec: FB2\} prints 7 and occurs at source lines 364 (twice), 541 and 927. In addition the article prints the label thm: spectral gap twice, once as a stray label on an otherwise empty line of the Introduction (source line 266) and once on the lemma it names (source line 787); the stray introduction occurrence is renamed, an invisible label that prints nothing, so that the delivered unit compiles without multiply-defined labels, and the references at source lines 264 and 794 continue to resolve to Lemma 4.2 at source line 787. Printed slips kept literal: no printed word, symbol, hypothesis or number of the counted statement or of its proof was altered, so the proof keeps its own typographic forms, including the restricted-norm notation in its last display, the doubled label of Theorem 1.2 in the source, and the source's habit of naming labels after other environments (the main theorem carries the labels thm: resolvent and thm: main, the spectral-gap lemma is labeled thm: spectral gap and the qualitative stability proposition is labeled p:standard); labels are invisible in print and are reproduced as printed. The source contains no non-ASCII character, no figure, no table, no \textbackslash{}includegraphics and no \textbackslash{}input anywhere in the e-print, so nothing had to be omitted for asset reasons and no artwork is redistributed. The wrapper loads mathalpha (the current name of the article's mathalfa) once with the pinned runtime options in place of the article's own scr=boondox and cal settings; because the wrapper's script family has no lowercase letters while the article's script alphabet does, the article's own script family is declared in the wrapper preamble exactly as mathalpha's scr=boondox option declares it in this runtime, so the script letters of the delivered unit (source lines 432, 448 and the three uses of the script-h macro) are set in the article's own Boondox calligraphic font and no glyph is missing. The article's epstopdf, color, url and graphicx loads and its .tif graphics rule are not needed by this unit and are not reproduced. No mathematics is written, repaired or re-derived; all mathematics is native text with no image dependency.
\section{Introduction}
Let  $\Omega \subset \R^n$ be an open bounded set. The Laplace equation with Dirichlet boundary data has  discrete spectrum $0 < \lambda_1(\W) <  \lambda_2(\W) \leq \dots$ of eigenvalues, whose corresponding eigenfunctions $\{ u_{\W, k} \}_{k \in \mathbb{N}}$ solve the eigenvalue problem
\[
\begin{cases}
    -\Delta u_{\W,k} = \lambda_k(\W) \, u_{\W,k} & \text {in } \W\\
  \quad \ \   u_{\W,k} = 0 & \text{on } \partial \W
\end{cases}
\]
and, after normalizing so that $\| u_{\W,k}\|_{L^2(\W)} = 1$ which we shall always do,  form an orthonormal basis for $L^2(\W)$. 

A classical shape optimization problem is the minimization of the first eigenvalue $\lambda_1(\W)$ among domains $\W$ of a fixed volume. Following Lord Rayleigh's 1894 conjecture, in the 1920s Faber \cite{Faber23} and Krahn \cite{Krahn25} independently showed that balls are the only solutions to this variational problem. This leads to the Faber-Krahn inequality:
\begin{equation}\label{eqn: FK intro}
    \lambda_1(\W ) \geq \lambda_1(B)  \quad \text{ if } \quad|\W| = \omega_n\,,
\end{equation}
with equality if and only if $\W = B(x)$ for some $x\in \R^n$ (up to capacity zero sets).
Here $B(x) = \{|y-x|<1\}$ is the unit ball centered at $x$,  $B= B(0)$, and $\om_n = |B|$. The proof of the Faber-Krahn inequality is based on rearrangement and the isoperimetric inequality and  is nowadays classical.

A fundamental question related to the Faber-Krahn inequality is that of stability: if a domain $\W$ with $|\W| = \omega_n$ {\it almost} achieves equality in \eqref{eqn: FK intro}, then is $\W$ {\it almost} a unit ball? The first ``almost'' is naturally measured by the deficit in \eqref{eqn: FK intro}: $\lambda_1(\W)  - \lambda_1(B)$. As for the second,  there are various ways in which one can quantify the proximity of $\W$ to the space of unit balls,  with the most suitable  notion dependent on context.

One classical geometric notion is the Fraenkel asymmetry, $\alpha(\W) = \inf_{x \in \R^n} |\W\Delta B(x)|$, measuring the volume of the symmetric difference between $\W$ and the nearest unit ball. Sharp quantitative stability for the Faber-Krahn inequality with respect to the Fraenkel asymmetry was shown  in the paper \cite{BDV15} of Brasco, De Philippis, and Velichkov;  see also \cite{Bhat01, FMP09} for earlier results.

A different way to quantify how close $\W$ is to the space of unit balls
is {\it spectrally}: Are the eigenvalues and eigenfunctions of $\W$ close to the corresponding eigenvalues and eigenfunctions of the nearest unit ball? 
This notion of distance and corresponding quantitative stability estimates  turn out to be particularly important in applications to harmonic analysis and free boundary problems.

These applications motivated the authors' previous work in \cite{AKN1, AKN2}, where it was  shown that for any open bounded set $\W \subset \R^n$ with $|\W|  =\omega_n$,
\begin{equation}\label{eqn: first efn stability}
\lambda_1(\W ) - \lambda_1(B) \geq c_n \inf_{x \in\R^n} \int |u_{\W,1 } - u_{B(x), 1}|^2 \,.
\end{equation}
Here the eigenfunctions are extended by zero to be defined on $\R^n$.
In the same work, an analogue of \eqref{eqn: first efn stability} was shown on the round sphere $\mathbb{S}^n$. This was used in \cite{AKN3} to establish a rectifiability and uniqueness-of-blowups theorem in terms of the Alt-Caffarelli-Friedman monotonicity formula, with applications to free boundary regularity for certain (e.g. vectorial) free boundary problems. In \cite{FTV1, FTV2}, Fleschler, Tolsa, and Villa applied the analogue of \eqref{eqn: first efn stability} on $\mathbb{S}^n$ as a main tool, alongside  some deep covering arguments, to establish a higher-dimensional version of Carleson's famed $\e^2$-conjecture in harmonic analysis \cite{JTV21}. In \cite{FTV1}, the authors also apply \eqref{eqn: first efn stability} to prove a different form of quantitative stability for the Faber-Krahn inequality with respect to a geometric distance involving the capacity.

The first main result of this paper substantially extends \eqref{eqn: first efn stability}: the Faber-Krahn deficit quantitatively controls not only the distance between first eigenfunctions,  but between  {\it all eigenfunctions}, of $\W$ and the nearest unit ball:


\begin{theorem}\label{thm: eigenfunction estimates} Fix $n \geq 2$ and $k \in \mathbb{N}$. There is a positive constant $c_{n,k}$ depending only on $n$ and $k$ such that 
for any open bounded set $\Omega \subset \R^n$ with $|\Omega| = \omega_n$, there is a point $x_\W \in \R^n$ such that
\begin{equation}\label{eqn: simple efn stability}
\lambda_1(\W ) - \lambda_1(B)  \geq c_{k, n} \int_{\mathbb{R}^n}\left|u_{\Omega, k}-u_{B({x_\W}),k}\right|^2\,.
 \end{equation}
\end{theorem}
As in \eqref{eqn: first efn stability} above, in the statement of Theorem~\ref{thm: eigenfunction estimates}, we extend each $u_{\W,k} \in H^1_0(\W)$ by zero to be in $H^1_0(\R^n)$, and likewise for eigenfunctions on the ball.  When $\lambda_k(B)$ has multiplicity greater than $1$, $u_{B({x_\W}),k}$ denotes {\it some} eigenfunction in the $\lambda_k(B)$-eigenspace with $L^2$ norm equal to $1$, since of course any choice of ordered basis for this eigenspace is not canonical. The quadratic power on the $L^2$-norm on the right-hand side of \eqref{eqn: simple efn stability} is optimal; see section~\href{https://arxiv.org/abs/2504.13053v2}{5.3}.

One can view Theorem~\ref{thm: eigenfunction estimates} as the counterpart for eigenfunctions of a recent result of Bucur, Lamboley, Nahon, and Prunier \cite{blnp23} for eigenvalues. There, the authors show  that
for each $\Omega \subset \R^n$ with $|\Omega| = \omega_n$ and  $k \in \mathbb{N}$, the following sharp quantitative estimates hold: if $\lambda_k(B)$ is simple, then 
\begin{align}\label{eqn: Bucur simple}
 \lambda_1(\W ) - \lambda_1(B)   & \geq c_{n,k}  |\lambda_k(\W)-\lambda_k(B)|,   \\
 \intertext{
while if $\lambda_k(B)$ has multiplicity---say $\lambda_{k-1}(B) <  \lambda_k(B)=\ldots=\lambda_{\ell}(B) < \lambda_{\ell+1}(B)$---then 
 }
 \label{eqn: Bucur multiplicity}
 \lambda_1(\W ) - \lambda_1(B)    &\geq c_{n,k}    \Big| \sum_{i=k}^\ell ( \lambda_i (\W) - \lambda_i(B) \Big| \,.
\end{align}
The linear powers in \eqref{eqn: Bucur simple} and \eqref{eqn: Bucur multiplicity} are stronger than one typically expects in stability estimates (and would be false for Theorem~\ref{thm: eigenfunction estimates}), and have to do with the fact that the ball is critical for the functionals on the right-hand sides, as well as the left-hand sides, of \eqref{eqn: Bucur simple} and \eqref{eqn: Bucur multiplicity}. 

Theorem~\ref{thm: eigenfunction estimates} is proven as an application of the other main theorem of the paper: sharp  quantitative estimates for the resolvent operator. To state the result we first introduce a bit of terminology. For an open bounded set $\W \subset \R^n$ and a function $f \in H^{-1}(\R^n)$, let $u_{\W}^f  \in H^{1}_0(\W)$ be the unique (weak) solution to the Poisson equation 
\begin{equation}\label{eqn: PDE with f RHS}
	\begin{cases}
	\hfill	-\Delta u^f_{\W}  = f & \text{ in } \W,\\
	\hfill	u_{\W}^f =0 & \text{ on } \partial \W.
	\end{cases}
\end{equation}
As above for eigenfunctions, we extend $u_\W^f$ by zero to be defined as a function in $H^1_0(\R^n)$. The {\it resolvent operator}
$$(-\Delta)_\W^{-1} : H^{-1}(\R^n) \to H^1_0(\R^n)$$ associates to $f \in H^{-1}(\R^n)$ the unique solution to \eqref{eqn: PDE with f RHS}. As an operator from $L^{\infty}$ to $L^2$, the resolvent operator of $\W$ is close to the resolvent operator of the nearest unit ball, quantitatively in terms of the Faber-Krahn deficit:
\begin{theorem}\label{thm: resolvent}\label{thm: main}
Fix $n\geq 2$. There is a constant $c_{n}>0$ depending on $n$ such that  for any bounded open set $\W \subset \R^n$ with $|\W| = \omega_n$,
\begin{align}
    \label{eqn: resolvent estimate}
  \lambda_1(\W ) - \lambda_1(B)  & \geq c_{n} \left\| (-\Delta)_{\W}^{-1} -(-\Delta)_{B(x_\W)}^{-1}\right\|^2_{L^{\infty}\to L^2}\,\\
\intertext{for a point $x_\W \in \R^n$. In other words, for any $f:\R^n \to \R$ with $\|f\|_{L^{\infty}(\R^n)}\leq 1,$ we have}
\label{eqn: quant stability main}
 \lambda_1(\W ) - \lambda_1(B)  & \geq c_{n}\int_{\R^n} \big|u_{\W}^f - u_{B(x_\W)}^f\big|^2 \,.   
\end{align}
\end{theorem}
To establish \eqref{eqn: resolvent estimate} (and in turn Theorem~\ref{thm: eigenfunction estimates}) from \eqref{eqn: quant stability main}, it is essential that the constant $c_{n}$ and the ball center  $x_\W$ in \eqref{eqn: quant stability main} do not depend on the function $f$. The point $x_\W$  is a truncated version of the barycenter of $\W$ defined  in section~\ref{sec: truncated barycenter}. 


A {\it  qualitative} version of \eqref{eqn: quant stability main}, with the modulus of continuity depending on $f$, can be deduced by combining a theorem of \v{S}ver\'{a}k  \cite[Theorem 3.2.5]{hp18}, the Kohler-Jobin inequality (see Lemma~\ref{lem: KJ consequence}), and a standard concentration compactness argument. To our knowledge, Theorem~\ref{thm: main} is the first {\it quantitative} estimate for the  resolvent operator.

 

The quadratic powers in \eqref{eqn: resolvent estimate} and \eqref{eqn: quant stability main} are optimal and cannot be replaced by a smaller exponent, as shown in section~\href{https://arxiv.org/abs/2504.13053v2}{5.3}. In this sense Theorem~\ref{thm: main} is sharp. A related question is whether in \eqref{eqn: resolvent estimate} the operator norm from $L^{\infty}$ to $L^2$ can be replaced with  the operator norm from $X$ to $Y$ for other function spaces $X\supset L^{\infty}$ and $Y \subset L^2$. We construct examples in section~\href{https://arxiv.org/abs/2504.13053v2}{5.3} exploring the extent to which one might expect to prove a ``strong form stability'' improvement in this direction. For instance, we show $L^2$ cannot be replaced by $H^1$, and $L^{\infty}$ cannot be replaced by $L^p_{loc}$ for $p<n/2.$ However, we have no examples ruling out the $X= L^p$ with $p>n$, which leaves the following natural open question:
\smallskip

\noindent{\bf Open Question.} {\it Does Theorem~\ref{thm: main} hold with the norm $L^p$ in place of the norm  $L^{\infty}$?
}

\smallskip
 We expect the answer to this open question to be affirmative. Many of the arguments work with $f \in L^p$ with $p>n$. 


As mentioned above, we also showed \eqref{eqn: first efn stability} on the round sphere and hyperbolic space in \cite{AKN1, AKN2}. Similarly, Theorems~\ref{thm: eigenfunction estimates} and \ref{thm: main} can be shown to hold on these spaces. Nearly  all of the arguments in this paper can be adapted to these settings, with the only nontrivial modifications being (a) the spectral gap of Lemma~\ref{thm: spectral gap} below should be replaced by a spectral gap argument by adapting the proof of \cite[Theorem 3.11]{AKN1} as well as the arguments in \cite{p24}, and (b) a minimizer of the penalized variational problem in Theorem~\ref{thm: summary for minimizers} below should be shown to satisfy the volume constraint by using the Euler-Lagrange equation as in \cite[Proposition 9.5]{AKN2} instead of the simpler scaling argument of Lemma~\href{https://arxiv.org/abs/2504.13053v2}{7.15}. To aid in the clarity of presentation, we have opted to present all the proofs here for domains on Euclidean space.

\label{thm: spectral gap (stray introduction occurrence)}




Beyond the works mentioned above, there is a rich literature on quantitative stability for geometric inequalities involving PDE-driven shape functionals like $\lambda_1$, with a number of exciting developments in the past decade or two. 
Following the work of Brasco-De Philippis-Velichkov \cite{BDV15}, similar methods were used to establish sharp quantitative stability for the $p$-Faber-Krahn \cite{FuscoZhang}, isocapacitary  \cite{DPMMCapacity}, and $p$-isocapacitary \cite{Mukoseeva+2021} inequalities.
Through quite different methods, quantitative stability has been established for spectral shape optimization problems with Neumann \cite{BrascoPratelli} and Robin \cite{BucurRobinLaplacian} boundary conditions, and on surfaces \cite{MR4859580} (both globally and within a fixed conformal class). The recent papers \cite{acampora2025sharpquantitativetalentisinequality, amato2024talenticomparisonresultquantitative} address stability for Talenti's rearrangement inequality. This sampling of results is by no means exhaustive, and we refer the reader to \cite{BrascoDePSurvey} for an overview of spectral inequalities in quantitative form.




\subsection*{Discussion of the Proofs}
Theorem~\ref{thm: eigenfunction estimates} will be established as a consequence of Theorem~\ref{thm: main}. The basic idea in the case when $k\in \mathbb{N}$ corresponds to a simple eigenvalue $\lambda_k(B)$ of the ball (and $x_\W = 0$ for simplicity) is as follows. We apply Theorem~\ref{thm: main} with a normalization of $f = \lambda_k(B) u_{B,k}$ to find $u_B^f = u_{B, k}$ is quantitatively close to $u_\W^f$. Then, expanding both functions in the basis of eigenfunctions on $\W$, we are able to quantitatively estimate their Fourier coefficients in order to show $u_{B,k}$ is quantitatively close to $u_{\W ,k}$.

Although the Faber-Krahn inequality is a consequence of the isoperimetric inequality, it is thus far unclear how to deduce sharp quantitative stability for the Faber-Krahn inequality (even with respect to the Fraenkel asymmetry) from quantitative stability for the isoperimetric inequality \cite{FuMaPr, FiMaPr, CiLe12}; see \cite{Melas, HansenNadir94, FMP09} for non-sharp or conditional results and the survey \cite{BrascoDePSurvey} for a discussion of the limitations of this approach.

Instead, we use the selection principle framework introduced by Cicalese and Leonardi in \cite{CiLe12}. This robust tool, which has found application in myriad settings, e.g. \cite{FuJu, BDF, BDS15, N16, DPMMCapacity, Mukoseeva+2021, FuscoZhang}, allows one to prove quantitative stability for a given geometric inequality from {\it linear stability} together with {\it regularity estimates} for a suitable penalized functional. 


Brasco-De Philippis-Velichkov's quantitative Faber-Krahn inequality \cite{BDV15} also used a selection principle scheme. While the linear stability portion of our proof follows in a straightforward manner from theirs, the penalized functional and corresponding regularity estimates that are needed to obtain the resolvent estimates \eqref{eqn: resolvent estimate}  differ fundamentally  from those needed to control the Fraenkel asymmetry in \cite{BDV15}. 
In this respect, a major new feature of this paper is the development of techniques introduced in our earlier paper \cite{AKN2} to study the regularity of minimizers of Alt-Caffarelli-type functionals involving a type of a-priori same-order perturbation described below.

First, several initial reductions simplify the proof of Theorem~\ref{thm: main} before running the selection principle. It suffices to show the estimate \eqref{eqn: quant stability main} in the case when
 \begin{equation}\label{eqn: f hyp after reductions}
     f \in C_0^\infty(\R^n), \qquad f\geq 0 , \qquad \|f\|_{L^{\infty}(\R^n)}\leq 1,
 \end{equation}
 as shown in section~\ref{sec: simplify}. Moreover, by the Kohler-Jobin inequality, it suffices to show Theorem~\ref{thm: main} with  the Faber-Krahn deficit $\lambda_1(\W) -\lambda_1(B)$ replaced by the Saint-Venant deficit $\tor(\W)-\tor(B)$. Here $\tor(\W)$ is
the torsional rigidity of $\W$,  defined as
	\begin{equation}\label{e: torsion def}
 \tor(\W) =\inf  \left\{ \int \frac{1}{2} |\grad w|^2 - w \, dx \ : \ w \in H^1_0(\W)\right\}.
	\end{equation}
The sign convention adopted here differs from the one used in some other references (and is minus the physical quantity). The torsional rigidity is  a non-positive quantity, and among sets of a fixed volume, it is uniquely (up to sets of capacity zero) minimized by balls. The resulting inequality is the {\it  Saint-Venant inequality}:
	\begin{equation}\label{e: saint venant}
    \tor(\W)\geq \tor(B) \qquad \text{ if } \quad|\W| = \omega_n\,.
    \end{equation}
The trick of replacing the Faber-Krahn deficit with the Saint-Venant deficit (and more generally, replacing $\lambda_1$ with the more-regularizing shape functional $\tor$) is well known and was used, for instance, in \cite{BDV15, blnp23}.

In view of these reductions, to prove Theorem~\ref{thm: main} it suffices to show there is a constant $c_{n}>0$ such that 
\begin{align}
\label{eqn: quant stability main 3}
\tor(\W ) - \tor_1(B)  & \geq c_{n}\int_{\R^n} \big|u_{\W}^f - u_{B(x_\W)}^f\big|^2 \,
\end{align}
for any open bounded set $\W \subset \R^n$ with $|\W| = \om_n$ and for any  $f:\R^n \to \R$ satisfying \eqref{eqn: f hyp after reductions}. Here, as in Theorems~\ref{thm: eigenfunction estimates} and \ref{thm: main} above, $x_\W$  denotes the truncated barycenter, defined in section~\ref{sec: truncated barycenter}. For sets of diameter at most $2$, like the minimizer in Theorem~\ref{thm: summary for minimizers} below, the truncated barycenter agrees with classical barycenter $\fint_\W x\,dx$. Of course, it is essential that  the constant $c_{n}$ in \eqref{eqn: quant stability main 3} does not depend on the smoothness of $f$ beyond its $L^{\infty}$ norm or on its compact support.

By using a selection principle reduction (carried out in section~\ref{sec: proofs of main theorems}), the main analytical challenge to prove \eqref{eqn: quant stability main 3} is to establish the existence and regularity of minimizers of the penalized functional 
\begin{equation}\label{eqn: energy functional intro}
		 \mathcal{E}(\W) = \tor(\W) + \tau \, 
        h\Big(\int |u_\W^f - u_{B(x_\W)}^f|^2  \Big)  
	\end{equation} 
among sets of volume $\om_n$. Here $h(t) =\sqrt{a^2 +(a - t )^2}$ for a given number $a\in(0,1]$ is (up to a constant) a regularization of the function $|t-a|$, and $\tau>0$ is a parameter that will be chosen to be small.
We show the following result for the volume-constrained minimization of the functional $\mathcal{E}$ in \eqref{eqn: energy functional intro}.


\begin{theorem}\label{thm: summary for minimizers} \label{summary ex and reg theorem}
Fix $n\geq 2$, $\alpha\in (0,1)$, and $\delta_0\in (0,1]$. There exists $\bar{\err}_0= \bar{\err}_0(n,\alpha, \delta_0)>0$
such that the following holds. Fix $a\in (0,1]$, $\tau \in (0,\bar{\tau}]$, and $f$ as in \eqref{eqn: f hyp after reductions}. The variational problem 
\begin{equation}
	\label{eqn: main min prob}
\inf \{\mathcal{E}(\W) : \W \subset \R^n \text{ open, bounded, } |\W|= \omega_n\}
\end{equation}
admits a minimizer $U$. There is a $C^{2, \alpha}$ function $\xi: \partial B(x_\W) \to \R$, with  $\| \xi \|_{C^{1,\alpha}(\partial B(x_\W))} \leq \delta_0$ such that 
 \[
 \partial U = \{ (1+\xi(x))x : x \in \partial B(x_\W) \}\,. 
 \]
\end{theorem}

{Since we assume in \eqref{eqn: f hyp after reductions} that $f$ is smooth, we obtain $\xi \in C^{2,\alpha}$. However, since we do not assume bounds on the derivative of $f$ we only obtain the first order bound $\|\xi\|_{C^{1,\alpha}}\leq \delta_0$, but utilizing the recent ideas in \cite{p24} to obtain Lemma \ref{thm: spectral gap} this will be sufficient for our methods.} 

Minimization of the first term in the energy $\mathcal{E}$, the torsional rigidity, behaves much like the minimization of the Alt-Caffarelli (also called one-phase Bernoulli) functional. Regularity theory for minimizers of this functional was first developed by Alt and Caffarelli in \cite{AC81} and its variants have since been a central topic in free boundary problems. 

What introduces new challenges to the analysis of the functional $\mathcal{E}$ is the term $\nl(\W):= h(\int |u_\W^f - u_{B(x_\W)}^f|^2)$, which is the same order as $\tor(\W)$, has poor lower semicontinuity properties, and in the Euler-Lagrange equation contributes a same-order term to the one coming from $\tor(\W)$ without a favorable sign.

Standard methods of proving existence fall short for the minimization problem \eqref{eqn: main min prob}; for instance 
\cite{BuDalMaso} cannot be applied as the functional $\mathcal{E}$ is not decreasing under set inclusions. Instead we follow the approach of our earlier paper \cite{AKN2}: we prove a priori estimates for inward and outward minimizers and construct a ``good'' minimizing sequence consisting of outward minimizers along which the functional actually is lower semicontinuous. This approach has the additional advantage that we work entirely in the framework of open sets and  there is no  need to pass to quasi-open sets.

In terms of regularity, since the term $\nl(\W)$ scales at the same order as $\tor(\W)$, one cannot apply the theory of almost-minimizers  of Alt-Caffarelli-type functionals \cite{DET}. A formal  computation  assuming $\partial \W$ is sufficiently regular shows that the Euler Lagrange equation is the vectorial Alt-Caffarelli problem
\begin{equation}
    \label{eqn: intro formal EL}
    \begin{cases}
\ \        -\Delta w_\W = 1, \ \  -\Delta u_\W^f = f, \ \ -\Delta p_\W = u_\W^f -u_{B(x_\W)}^f & \text{ on } \Omega\\
    \ \  \ \ \ \  \  w_\W = u_\W^f = p_\W = 0 & \text{ on } \partial \W,\\
       \ \  -|\na w_\W|^2 + C_\W \tau ( |\nabla p_\W| \, |\nabla u_\W^f| - g_\W)  = \text{const} & \text{ on }\partial \W\,.
    \end{cases}
\end{equation}
Here $g_\W:\R^n \to \R$ is a smooth bounded function depending on $\W$ (through the solution of another auxiliary PDE) and  $C_\W = h'(\int |u_\W^f - u_{B(x_\W)}^f|^2)>0$.
Significant difficulties arise in rigorously deriving the stationarity condition \eqref{eqn: intro formal EL}. Stemming from the genuinely nonlocal nature of the functional $\nl(\W)$---which leads, in particular, to the presence of the auxiliary potential $p_\W$ in \eqref{eqn: intro formal EL} that solves a Poisson equation with right-hand side depending on $\W$---it is not clear how to run the type of blowup argument used in \cite{blnp23} to show $\W$ is a viscosity solution of \eqref{eqn: intro formal EL}. 

Instead, we use careful regularization arguments to establish a distributional form of \eqref{eqn: intro formal EL} and then derive a pointwise form by plugging in suitable test functions (with the final line of \eqref{eqn: intro formal EL} understood in the sense of nontangential limits). On a broad level, this approach follows the one of our earlier paper \cite{AKN2}, but we are able to refine a number of our earlier arguments, leading to a significantly simpler proof.


Once we are able to rigorously establish that \eqref{eqn: intro formal EL} holds (and along the way, prove the $\W$ is an NTA domain), we can apply the inhomogeneous boundary Harnack theorem \cite{aks23} shown by the first two authors and Shahgholian. This allows us to recast \eqref{eqn: intro formal EL} as a standard one phase problem and apply the classical theory of Alt and Caffarelli \cite{AC81}.


\subsection*{Organization of the Paper}

The paper  essentially splits into two disjoint parts. After recalling notation and classical facts from the literature in section~\ref{sec: prelim},  we dedicate the first part of the paper, sections~\ref{sec: more preliminaries}--\ref{sec: proofs of main theorems}, to proving Theorems~\ref{thm: eigenfunction estimates} and \ref{thm: main} assuming Theorem~\ref{thm: summary for minimizers}. Specifically, we introduce the truncated barycenter and establish several reductions for our main theorems in section~\ref{sec: more preliminaries}, show  Theorem~\ref{thm: main} in the class of nearly spherical sets in section~\ref{sec: linear stability}, and in section~\ref{sec: proofs of main theorems} we conclude the proofs of Theorems~\ref{thm: eigenfunction estimates} and \ref{thm: main} and present examples demonstrating the sharpness of the theorems. 

The second part of the paper, sections~\href{https://arxiv.org/abs/2504.13053v2}{6} and \href{https://arxiv.org/abs/2504.13053v2}{7},  is dedicated to proving Theorem~\ref{thm: summary for minimizers} using techniques from  free boundary theory. In section~\href{https://arxiv.org/abs/2504.13053v2}{6} we establish the existence of minimizers of a relaxed version of the minimization problem \eqref{eqn: main min prob}, as well as basic properties of minimizers like Lipschitz regularity and nondegeneracy of the torsion function. In section~\href{https://arxiv.org/abs/2504.13053v2}{7}, we compute the Euler-Lagrange equation satisfied by minimizers and use this to show boundary regularity and that the minimizer of the relaxed problem satisfies the volume constraint $|U| =\om_n$, making it a minimizer of the original minimization problem \eqref{eqn: main min prob}.  
% 
\medskip


\noindent{\it Acknowledgments.} MA is supported by Simons award 637757.
DK is supported by NSF grant DMS-2247096.
RN is supported by NSF grants DMS-2340195, DMS-2155054, and DMS-2342349. The authors are very grateful to Jimmy Lamboley for showing us Prunier's paper \cite{p24}; adapting an argument therein allowed us to substantially weaken the hypotheses of our main theorems.

\section{Preliminaries}\label{sec: prelim}
We introduce some known preliminary facts and notation that will be used throughout the paper. For the remainder of the paper, we fix $n$ to be an integer with $n\geq 2.$
 We write $B_r(x)$ to denote a ball of radius $r$ centered at $x$, and when $r=1$ or $x=0$ we often omit the dependence in the notation, writing, e.g. $B_r = B_r(0)$ and $B(x) = B_1(x).$


\subsection{Basic facts about the torsional rigidity}
Let $\W\subset \R^n$ be an open bounded domain. The infimum in the variational problem \eqref{e: torsion def} is uniquely achieved by the {\it torsion function} $w_\W \in H^1_0(\W)$ solving
\begin{equation}
	\label{e: torsion eqn}
\begin{cases}
\hfill	-\Delta w_\W = 1 & \text{ in }\W\\
\hfill	w_\W = 0 & \text{ on } \partial \W.
\end{cases}
\end{equation}
Of course, \eqref{e: torsion eqn} is a special case of  \eqref{eqn: PDE with f RHS} with $f \equiv 1$, so that $w_\W = u_\W^1 $ in the notation introduced there. As for eigenfunctions and solutions to \eqref{eqn: PDE with f RHS}, we extend $w_\W$ by zero to be defined in $H^1_0(\R^n)$, using the same notation $w_\W$ to denote the extended function. The distributional Laplacian on $w_\W$ (defined by acting on $\phi \in C^\infty_c(\R^n)$ by $\Delta w_\W(\phi) = \int w_\W \Delta \phi$) satisfies
\begin{equation*}
    -\Delta w_\W \leq 1 \quad\text{ on } \R^n\, .
\end{equation*}
The same is true for any solution of \eqref{eqn: PDE with f RHS} with $\| f\|_{L^\infty(\R^n)}\leq 1.$
Multiplying \eqref{e: torsion eqn} by $w_\W$ and integrating over $\R^n$, we obtain the following expressions for the torsional rigidity:
\begin{equation}\label{e: torsion and dir energy}
\tor(\W) = -\frac{1}{2} \int |\na w_\W|^2 = -\frac{1}{2} \int w_\W.
\end{equation}
If $\W'\subset \W,$ then any competitor $w $ for $\tor(\W')$ is also a competitor for $\tor (\W)$ and thus $\tor(\W) \leq \tor(\W')$. Direct computation shows $\tor\left(r \W \right) = r^{n+2} \tor (\Omega)$ and that 
the torsion function on a ball is 
\begin{equation}
 	\label{eqn: tor function on a ball}
 w_{B_r(x_0)} = \frac{1}{2n}(r^2 - |x-x_0|^2 ) \chi_{B_r(x_0)}. 
  \end{equation}
In particular, $-\tor(B_1)  = \frac{\omega_n}{2n(n+2)}$ and $\| w_{B_r(x_0)}\|_{L^\infty(\R^n)} =\frac{r^2}{2n}$.

By the maximum principle, $\| u_\W^f\|_{L^\infty}  \leq \| w_\W\|_{L^\infty}$, while
 Talenti's theorem\footnote{if $u=u_\W^f$ solves \eqref{eqn: PDE with f RHS} 
and $v= u_{B_r}^{f^*}$ solves \eqref{eqn: PDE with f RHS} on $B_r$ 
 where $r>0$ is again chosen to that $|\W| = |B_r|$ and $f^*$ is the symmetric decreasing rearrangement of $f$, then $u^* \leq v$ pointwise on $B_r$. } \cite{Talenti1976} guarantees that $\| w_\W\|_{L^\infty} \leq \|w_{B_r}\|_{L^\infty}$, where $r$ is such that $|B_r|=|\W|$. 
Thus, for any bounded open set $\W \subset \R^n$ with $|\W| \leq 2 \omega_n$, 
\begin{equation}
    \label{eqn: bound for tor fn}
\| w_\W\|_{L^\infty(\R^n)} \leq \frac{2}{n}, \qquad \| w_\W\|_{L^2(\R^n)}^2 \leq \frac{4\om_n}{n^2} \leq 1
\end{equation}


The Kohler-Jobin inequality \cite{KJ1, KJ2, KJ3}
says that among sets $\Omega \subset \mathbb{R}^n$ of a fixed torsional rigidity $\operatorname{tor}(\Omega)$, balls minimize the first Dirichlet eigenvalue $\lambda_1(\Omega)$. 
A well-known consequence of the Kohler Jobin inequality (see \cite[Cor. 3.18]{AKN1} or \cite[Prop. 2.1]{BDV15}) is that the deficit in the Faber-Krahn inequality linearly controls the definit in the Saint-Venant inequality:
\begin{lemma}\label{lem: KJ consequence}
 There exists a constant $C_n$ so that for any open set $\W\subset \R^n$ with $|\W|=\om_n$, 
\[
C_n\left(\lambda_1(\Omega)-\lambda_1(B)\right) \geq \operatorname{tor}(\Omega)-\tor(B)
\]
\end{lemma}


%
%
\subsection{Qualitative stability and relaxing the volume constraint}\label{ssec: qual stab}
As is typical in geometric variational problems, it is easier to replace a volume constraint with a volume penalization term. 
Following \cite{aac86} and \cite{BDV15}, we remove the volume constraint in \eqref{eqn: main min prob} in favor of adding to the energy a volume penalization term
\begin{equation} \label{eqn: def vol penalization fixed param}
	\mathscr{V}(t)=\fv(t) = \begin{cases}
		\vpar (t - \om_n) & t \leq \om_n\\
		\frac{1}{\vpar}(t - \om_n) & t > \om_n\,,
	\end{cases}
\end{equation}
where $\vpar>0$ is a parameter to be fixed later on.
By choosing  $\vpar$ sufficiently small, a set $\W$ whose volume exceeds the volume $\om_n$ of a unit ball will have a large energy contribution coming from the term $\fv(|\W|)$, and minimizing $\tor(\cdot)$ among sets of volume $\w_n$ is equivalent to the unconstrained problem of minimizing $\tor(\cdot)  + \fv(|\cdot|)$ among open bounded sets of volume at most $2\omega_n$:
\begin{lemma}
    \label{prop: base summary}
 There exists $\vpar_0 = \vpar_0(n)$ such that for $\vpar \leq \vpar_0$,  unit balls are the unique minimizers of the energy $\tor(\W)  + \fv(|\W|)$ among all open bounded sets $ \W \subset \R^n$ with $|\W | \leq 2 \omega_n$.
 Moreover, if $\W$ is an open bounded set with $|\W | \leq 2 \omega_n$ and  $\tor(\W)  + \fv(|\W|) \leq \frac{5}{4}\tor(B_1)$, then $|\W| \geq \omega_n/2$.
\end{lemma} 
The proof of Lemma~\ref{prop: base summary} is elementary and nearly identical to the proof of \cite[Lemma 4.5]{BDV15}, so we omit it.
%
The assumed hard volume constraint $|\W| \leq2 \om_n$ in Lemma~\ref{prop: base summary} is necessary, since for any fixed $\eta >0$,  a scaling computation shows that  $\tor(B_R) + \fv(|\W|) \to -\infty$ as $R \to \infty$ .  An alternative remedy, adopted in \cite{BDV15}, would be to minimize among sets contained in $B_R$, but this has the drawback of introducing a dependence of $\eta$ on $R$.
 
Next, a standard concentration compactness argument establishes the following qualitative stability result for the energy $\tor(\cdot) +  \mathscr{V}_\eta(|\cdot|)$: as $\tor(\W) +  \mathscr{V}_\eta(|\W|)$ tends to its infimum $\tor(B_1)$, $\Omega$ converges in $L^1$ to a ball:
 \begin{proposition}[Qualitative stability]
 	\label{lemma: qual stability}
 	Fix $\e>0$. There exist $\eta_1={\eta}_1(n)\in (0,\eta_0]$ and $\ccdel = \ccdel(n,\e)>0$ such that the following holds. Suppose  $\eta \in (0, \eta_1]$ and $\W\subset\R^n$ is an open bounded set with $|\W|\leq 2\omega_n$ and 
 	\begin{equation}
 		\label{eqn: small base energy}
 	\tor(\W) +\fv(|\W|) \leq \tor(B_1) + \ccdel.
 	 	\end{equation}
 	There exists $x_1$ such that
 	\begin{equation}
 		\label{eqn: qual stab}
 	|\W \Delta B(x_1) | \leq \frac{\omega_n\e}{2}.
 		\end{equation}
 \end{proposition}



%

\subsection{Geometric measure theoretic preliminaries} Let us recall some geometric measure theoretic definitions and facts that will be used to establish initial regularity and to compute the Euler-Lagrange equation of minimizers of the (relaxed version of the) minimization problem \eqref{eqn: main min prob}.



For $K>0$,
A bounded open set $\W\subset \R^n$ is called a {\it non-tangentially accessible (NTA) domain}  with NTA constant $K$ if the following properties hold: 
\begin{enumerate}
    \item  For each $x \in \partial \W$ and $r \in (0,1)$ there are points $y_{\textrm{i}}$ and $y_{\textrm{o}}$ such that 
    \[
    B_{r/K}(y_{\textrm{i}}) \subset B_r(x) \cap \W  \qquad \text{ and } \qquad B_{r/K}(y_{\textrm{o}}) \subset B_r(x) \setminus \W .
    \]
    \item  For each  $x,y \in  \W$, there is a curve $\gamma : [0,1] \to \W$ with $\gamma(0)=x$ and $\gamma(1)=y$ such that $\ell(\gamma([0,1]) \leq K |x-y|$ and 
    $$\text{dist}(\gamma(t),\partial \W) \geq \frac{1}{K} \min\{ \ell(\gamma([0,t])), \ell(\gamma([t,1]).
    $$
\end{enumerate}
\begin{remark}
    {\rm 
    Jerison-Kenig's original definition of NTA domain  \cite{JKNTA} takes a slightly different form but is equivalent to the one given here; see \cite[Theorem 2.15]{ahmnt17}.
    %for the proof of one implication of this equivalence (which the authors describe as a folklore theorem) and the references therein for the opposite implication. 
    Some references, e.g., the classical texts \cite{k94, CaffSalsa}, define NTA domains with the exterior corkscrew condition $B_{r/K}(y_{\textrm{o}}) \subset B_r(x) \setminus \W$ replaced by the weaker condition that $\Omega$ satisfies uniform upper volume density estimates at each $x\in \partial \W$. Since this definition is strictly weaker, the results cited below from \cite{k94}  hold for NTA domains as defined above. }
\end{remark}


For   a open bounded set $\W \subset \R^n$,  $x \in \partial \W$, and $\beta>0$, define the non-tangential approach region
\[
\Gamma_\beta(x)  = \{ y \in \W : |y-x| \leq (1+\beta) \text{dist}(y, \partial \W)\}.
\]
Given $u \in C(\Omega)$, the associated non-tangential maximal function $u^*:\pa \W \to \R$ (which depends on $\beta$) is defined by $u^*(y) = \sup_{x \in \Gamma_\beta(y)}|u(x)|$. The function $u$  is said to {\it converge non-tangentially to $a \in \R$} at $x \in \partial \Omega$ if, for every $\beta >0$, 
\begin{equation}\label{def: NTlimit 1}
\lim_{\substack{y \to x ,\\ 
y \in \Gamma_\beta(x)  }} u(y) = a\,.
\end{equation}

The next lemma gives us a Poission kernel representation for certain solutions of the Dirichlet problem when $\W$ is sufficiently regular. The statement is not the most general one possible, but will be sufficient for our needs.
\begin{lemma}\label{lem: poisson kernel}
    Let $\W \subset \R^n$ be an NTA domain and assume there are  constants $c,r_0 >0$ such that for each $x \in \partial \W$ and $r\in (0,r_0]$,
\begin{equation}\label{eqn: prelim ahlfors}
   c \, r^{n-1} \leq   \mathcal{H}^{n-1} (\partial \W \cap B_r(x) ) \leq \frac{1}{c} r^{n-1}\,.
\end{equation}
Then for each $x \in \W$, there is a function $K(x,y)$ defined  for $\mathcal{H}^{n-1}$-a.e. $y \in \partial \W$ such that the following holds. 
%
For $g\in L^\infty(\partial \W,  \mathcal{H}^{n-1})$, there is a unique harmonic function $u \in C^\infty(\W)$ with $u^* \in L^\infty(\partial \W, d\mathcal{H}^{n-1})$ for some $\beta>0$ such that $u$ converges non-tangentially to $g$ for $\mathcal{H}^{n-1}$-a.e. $y \in \partial \W$. For each $x \in \W$, $u(x)$ is given by
$$
u(x) = \int_{\partial \W} g(y){K}(x,y) \, d \mathcal{H}^{n-1}(y)\,.
$$
\end{lemma}
\begin{proof}
Suppose $\W$ is an NTA domain. Let $\{\w^x\}_{x\in \W}$ be the harmonic measure for $\W$, and for a fixed point $x_0 \in \W$, set $\w= \w^{x_0}$. By [Lemma 1.2.7]\cite{k94}, the measures $\w^x$ and $\w$ are mutually absolutely continuous for any $x \in \W$, so 
the Radon-Nykodym  derivative $\tilde{K}(x,y)$ of $\w^x$ with respect to $\w$
 exists for $\w$-a.e. $y$. By \cite[Corollary 1.4.3, Theorem 1.4.4]{k94}, for any $g \in  L^\infty(\pa \W, d \om)$, the function 
$u(x) = \int_{\partial \W} g(y)\tilde{K}(x,y) \, d\om(y)$
is the unique harmonic function on $\W$ with $u^* \in L^\infty(\partial \W, d\om)$ such that  $u(x)$ converges non-tangentially to $g$ for $\om$-a.e. $y \in \partial \W$. 

Now additionally assume that $\W$ satisfies \eqref{eqn: prelim ahlfors}. By \cite[Theorem 2]{DavidJerison}, the harmonic measure $\w$ and surface measure $\sigma:= \mathcal{H}^{n-1}\llcorner \partial \W$ are mutually absolutely continuous (in fact, quantitatively so). Thus, in each statement above, we may replace $\om$ by $\sigma$. This proves the lemma.
\end{proof}


A Borel set $\W\subset \R^n$ is a {\it set of finite perimeter} in $\R^n$ if 
\[\sup\Big\{ \int_\W \text{div}\,  T  : T\in C^1_c(\R^n; \R^n) , |T|\leq 1\Big\}<\infty.\]
If the topological boundary $\partial \W$ of measurable set $\W$ has  $\mathcal{H}^{n-1}(\partial \Omega) <\infty$, then $\W$ is a set of finite perimeter.
The reduced boundary $\partial^*\W \subset \partial \W$ of a set of finite perimeter $\W$ is the set of points such that 
\[
\nu_\W(x) := \lim_{r\to 0^+} \frac{D\chi_\W(B_r(x))}{|D\chi_\W| (B_r(x))} \qquad \text{ exists and has modulus }1.
\]
For every $x\in \partial^*\W$, the rescaled sets $\W_{x,r}:=(\W-x)/r$ converge in $L^1_{loc}$ to the half-space $H_x:= \{y \in \R^n : y \cdot \nu_\W(x)\} \leq 0$ as $r \to 0^+$; see \cite[Theorem 15.5]{MaggiBook}.
 If additionally there are positive constants $c, r_0>0$ such that
\[
  c\, \om_nr^n  \leq |B_r(x) \cap \W| \leq (1-c)\om_n r^n
\]
  for all $x \in \partial \W$ and $r\in (0,r_0]$, then $\mathcal{H}^{n-1}(\partial \W \setminus \partial^*\W) =0$ by \cite[Theorem 16.2]{MaggiBook}, and for any $x \in \partial^*\W$ and $\e>0$, we have $d_H(\partial \W \cap B_r(x) , \partial H_{x} \cap B_r(x)) < \e r$ for $r$ sufficiently small, and thus for $r$ small enough,
\begin{equation}\label{eqn: cone}
    \partial\W \cap B_r(x)\subset \{ y \in B_r(x): |y\cdot \nu_\W(x)| <\e|y|\} 
\end{equation}

The divergence theorem holds for vector fields $Z \in C^1_c(\R^n ; \R^n)$ on a set of finite perimeter $\W$. In section~\href{https://arxiv.org/abs/2504.13053v2}{7}, we will need to apply a different form of the divergence theorem, which holds assuming more regularity on $\W$ but less regularity for the vector fields. The following is a special case of \cite[Theorem 2.3.1]{HMT}. 
\begin{lemma}\label{l:fancy div thm}
Let $\Omega \subset \mathbb{R}^{n}$ be a bounded open set of finite perimeter with $\mathcal{H}^{n-1}(\partial \W \setminus \partial^*\W)=0$ and assume there are constants $c, r_0>0 $ such that
\eqref{eqn: prelim ahlfors} holds
for all  $r \in (0,r_0]$ and $x \in \partial \W$.
Then
\[
\int_{\Omega} \operatorname{div} Z d x =\int_{\partial \Omega} Z \cdot \nu_\W \,  d\mathcal{H}^{n-1}
\]
for any vector field $Z \in C(\Omega)$ with $\operatorname{div} v \in L^1(\Omega)$ and $Z^* \in L^\infty( \partial \W, \mathcal{H}^{n-1})$ whose pointwise nontangential limit exists $\mathcal{H}^{n-1}\llcorner \partial \W$-a.e.
\end{lemma}
%%
%
%
%



\section{The truncated barycenter and initial simplifications}\label{sec: more preliminaries}
We define the truncated barycenter and prove a key Lipschitz property in Lemma~\ref{lem: lip bary}, then prove three lemmas that reduce the proof of  Theorem~\ref{thm: main} to $f$ that are smooth, nonnegative, and compactly supported.

 \subsection{Truncated barycenters and a key Lipschitz property}\label{sec: truncated barycenter}
\label{ssec: lip bary}
%
%
%
%
Simple examples, for instance, the sequence of sets $\Omega_j = B_{1-1/j}(0) \cup B_{r_j}(2^j\,e_1)$ with $r_j$  chosen so $|\Omega_j|= \om_n$,  show that even a qualitative form of Theorem~\ref{thm: main} is false if $x_\W$ is chosen to be the barycenter $ \tfrac{1}{|\W|}\int_\W x \,dx$ of $\W$.
This is a familiar challenge in stability problems, and the usual solution is to take the infimum over all $x \in \R^n$ of the given distance $d(\W, B_1(x))$, rather than choosing an explicit ball center. This solution is unsuitable in the present context for two reasons: first, it leads to serious challenges in establishing the necessary regularity for minimizers. 
Second, to prove Theorem~\ref{thm: eigenfunction estimates}, it is essential that the ball center in \eqref{eqn: quant stability main} is independent of the function $f$.\footnote{The first of these two issues already arises in \cite{BDV15}. 
Their solution is to prove a quantitative stability for sets contained in  $B_R$ using the barycenter $\text{bar}(\W)$, then pass $R$ to infinity using a previously-known non-sharp quantitative stability estimate \cite{FMP09} for the Faber-Krahn inequality with respect to the asymmetry. No analogous non-sharp result exists in the present context.}


In view of this, we define a suitable {truncated barycenter} $x_\W$ that agrees with $\text{bar}(\W)$ for sufficiently well-behaved sets but de-emphasizes small pieces of $\W$ at infinity, so that, for instance, $x_{\Omega_j} \to 0$ for the sequence $\W_{j}$ above. In view of the characterization  of the classical barycenter as $ \text{argmin}\{ \int_\Omega |x-y|^2 \,dy : x \in \R^n\}$, we define the {\it truncated barycenter} $x_\W$ of a measurable set $\W$ with $|\W|>0$ by 
\begin{equation}
	\label{eqn: barycenter definition}
x_\W = \text{argmin}\Big\{ \int_\W \TBF(|x-y|)\,dy : x \in\R^n \Big\},
\end{equation}
where  $\TBF: \R_+\to \R_+$ is a fixed $C^2$ function satisfying 
\begin{equation}
    \label{eqn: conditions on q}
\TBF(t) = t^2\text{ for }t\leq 100, 
\qquad \TBF'(t)<250 \text{ for }t\in \R, \quad \text{ and } \quad  0 < \TBF'' (t)<3 \text{ for }t\in \R.
\end{equation}
A simple way to construct such a function is to let $\ddot{q}(\rho)=2$ for $\rho\leq 100$ and $\ddot{q}(\rho) =\tfrac{1}{\rho^2 -100^2+1/2}$ for $\rho \geq 100$ and take $q(t) = \int_0^t\int_0^s \ddot{q}(\rho)\,d\rho\,ds$.  
It is easy to check from \eqref{eqn: conditions on q} that $x\mapsto \int_\W \TBF(|x-y|)\,dy$ is a strictly convex function and thus $x_\W$ is well-defined, and that $x_\W$ agrees with the classical barycenter for those sets $\W$ with diameter at most $100$.
An important feature of the truncated barycenter is the following Lipschitz continuity property:
 %
 %
 %
 \begin{lemma}\label{lem: lip bary}
  There are dimensional constants $\baryEps >0$ and $C>0$ such that the following holds. Let $\W, \W'\subset \R^n$ be bounded open sets with $|\W \Delta B| \leq \baryEps \omega_n$ and $|\W' \Delta B| \leq \baryEps \w_n$.  Then $x_\W, x_{\W'} \in B_3$ and  
  \begin{equation}
      \label{eqn: bary lip}
  |x_\W -x_{\W'}|\leq C |\W\Delta \W'|.
   \end{equation}
  Moreover, for any $f \in L^\infty(\R^n)$  with $\|f\|_{L^\infty(\R^n)} \leq1$,
\begin{equation}\label{eqn: beta lip}
\Big| \int   |u_\W^f - u_{B(x_\W)}^f|^2 -  \int  |u_{\W'}^f-u_{B(x_{\W'})}^f|^2   \Big| \leq C\,\Big( |\W\Delta \W'| +\int |u_{\W}^f - u_{\W'}^f|\Big)\,. 
\end{equation}

\end{lemma}


\begin{proof}
{\it Step 1:}	Let $\baryEps=\baryEps\in(0,1/5)$ be a fixed constant to be specified later in the proof. We first show $x_\W, x_\W' \subset B_3$. Taking the origin as a competitor in the minimization problem \eqref{eqn: barycenter definition} defining $x_\W$ shows 
	\[
	\int_{\W \cap B} \TBF(|x_\W-y|)\,dy \leq \int_{\W \cap B} \TBF(|y|)\,dy +  \int_{\W \setminus B} [ \TBF(|y|) - \TBF(|x_\W-y|) ] \,dy\leq \w_n + 250 \baryEps \omega_n |x_\W|.
	\]
Since $\TBF(t) \geq \TBF(1)+ \TBF'(1)(t-1) = 2t-1$ by convexity, the left-hand side has the lower bound
\[
\int_{\W \cap B} \TBF(|x_\W-y|)\,dy \geq 2|x_\W|(1-\baryEps) \omega_n - \int_{B} 2|y|+1 \, dy \geq 2|x_\W|(1-\baryEps) \omega_n - 3\omega_n\,.
\]
Choosing $\baryEps$ small enough, these two inequalities together show that $|x_\W| \leq 4/(2(1-\baryEps) -250\baryEps) \leq 3.$\\

{\it Step 2: } Next we show \eqref{eqn: bary lip}. The Euler Lagrange equation for $x_\W$ is 
$0 = \int_\W \TBF'(|x_\W-y|) \frac{x_\W - y}{|x_\W-y|} \,dy,$ so splitting the domain of integration into $\W \cap B_{50}$ and its complement yields the following expression for $x_\W$:
\[
x_\W = \frac{1}{|\W \cap B_{50}|}\Big( \int_{\W\cap B_{50}} y\,dy + \int_{\W \setminus B_{50}} \TBF'(|x_\W -y|) \tfrac{ x_\W -y}{|x_\W -y|} \,dy \Big)
\]
So, $
|x_\W -x_{\W'}| \leq I^{(1)} + I^{(2)} + I^{(3)}$
where 
\begin{align*}
	I^{(1)} &:=  \left| \frac{1}{|\W \cap B_{50}|} -  \frac{1}{|\W' \cap B_{50}|}\right| \, 
	 \bigg| \int_{\W\cap B_{50}} y\,dy + \int_{\W \setminus B_{50}} \TBF'(|x_\W -y|) \tfrac{ x_\W -y}{|x_\W -y|} \,dy \bigg|\\
	 I^{(2)} & : = \frac{1}{|\W' \cap B_{50}|}
	 \left|  \int_{\W' \cap B_{50}} \Big( y + \TBF'(|x_\W - y|) \tfrac{x_\W - y}{|x_{\W }-y|} \Big) \,dy -   \int_{\W \cap B_{50}} \Big( y + \TBF'(|x_\W - y|) \tfrac{x_\W - y}{|x_{\W }-y|} \Big) \,dy\right|\\
	 I^{(3)} &: =\frac{1}{|\W' \cap B_{50} |} \left| \int_{\W' \setminus B_{50}} \Big(\TBF'(|x_\W - y|) \tfrac{x_\W - y}{|x_{\W }-y|}  -\TBF'(|x_{\W'} - y|) \tfrac{x_{\W'} - y}{|x_{\W' }-y|} \Big)\,dy  \right|  \,.
\end{align*}
We have 
\begin{align*}
	I^{(1)} 
	 &\leq  \left( 51  \omega_n +250 \baryEps \omega_n\right) \left| \frac{1}{|\W \cap B_{50}|} -  \frac{1}{|\W' \cap B_{50}|}\right| 
	  \leq 200\omega_n^{-1} \big| |\W \cap B_{50}| - |\W' \cap B_{50}|\big| \leq 200\omega_n^{-1}|\W \Delta \W'|.
\end{align*}
Moreover, the modulus of the integrand in $I^{(2)}$ is bounded by $300$ and thus  
\[
I^{(2)}  \leq \frac{300 |\W\Delta \W'| }{|\W' \cap B_{50}|} \leq 600 \omega_n  |\W\Delta \W'|.
\]
 For the final term $I^{(3)}$, direct computation shows that  the function ${g}_y(x):= \TBF'(|x-y|) \frac{x-y}{|x-y|}$ has $|\nabla_x {g}_y(x)|\in C_n$ uniformly bounded for all $x,y$ with $|x-y| \geq 10$. In particular, since $x_\W, x_{\W'}$ lie in $B_3$, by the fundamental theorem of calculus,
\begin{align*}
		I^{(3)} \leq \frac{1}{|\W' \cap B_{50}|} 
		 \int_{\W' \setminus B_{50}} 
		 |{g}_y(x_\W) - {g}_y(x_\W')| \leq C_n \baryEps  |x_\W - x_\W'| \,. 
\end{align*}
Putting these estimates together we have $|x_\W -x_{\W'}| \leq C_n ( |\W\Delta \W'| + \baryEps |x_\W -x_{\W'}|)$. We conclude the proof of \eqref{eqn: bary lip} by choosing $\baryEps$ small enough to absorb the second term into the left-hand side.\\


{\it Step 3:} Now we prove \eqref{eqn: beta lip}. Fix $f \in L^\infty$ with $\|f \|_{L^\infty(\R^n)}\leq 1$. For brevity we write $u_\W, x,$ and $x'$ in place of $u_{\W}^f, x_{\W},$ and $ x_{\W'}$ respectively. 
 Since $u_\W, u_{\W'}, u_{B(x)}$ and $u_{B(x')}$  are each pointwise bounded above by $C_n$, we use the identity $a^2-b^2 =(a+b)(a-b)$ to find
	\begin{equation}\label{eqn: beta first}
    \Big|  \int \left( |u_\W - u_{B(x)}|^2 - |u_{\W'} -u_{B(x')}|^2   \right) \Big|
		 \leq C_n \Big(  \int\left|u_\W - u_{\W'} \right| + \int \left| u_{B(x)} - u_{B(x')}\right| \Big).
	\end{equation}
 By the triangle inequality we have 
	\begin{equation}
		\label{e: step 1}
		\begin{split}
	\int |u_{B(x)} - u_{B(x')}| 
	&= \int_{B(x) \setminus B(x')} 	|u_{B(x)}|  + \int_{B(x') \setminus B(x)} 	|u_{B(x')} | \\
	&+ \int_{B(x)\cap B(x')} |u_{B(x)} - u_{B(x')}|  
	 =: I^{(1)} + I^{(2)} +I^{(3)} 
	 		\end{split}
		\end{equation}
By the maximum principle and the expression \eqref{eqn: tor function on a ball} for the torsion function on a ball,  \begin{equation}\label{eqn: I1 and I2}
I^{(1)} + I^{(2)} \leq C_n |B(x) \Delta B(x')| \leq C_n  |x-x'|.
\end{equation}
Next, $v= u_{B(x)} - u_{B(x')}$ is harmonic on $E:=B(x)\cap B(x')$ so by the maximum principle and \eqref{eqn: tor function on a ball},
\[
\sup_{E} |v|= \sup_{ \partial E}|v| \leq \sup_{ \partial B(x)} |u_{B(x')}| + \sup_{  \partial B(x')} |u_{B(x)}| \leq 2 \sup_{\partial B(x)} w_{B(x')} \leq C_n |x-x'|.
\]
Integrating this bound over $E$ shows $I^{(3)} \leq C_n  |x-x'|.$
Combining this estimate with \eqref{eqn: beta first}, \eqref{e: step 1}, and \eqref{eqn: I1 and I2}, we see that the left-hand side of \eqref{eqn: beta first} is bounded above by $ C( \int |u_\W - u_{\W'}| + |x - x'|).$
Finally,  choosing $\baryEps$  according to Lemma~\ref{lem: lip bary} so that $
|x-x'| \leq C_n |\W \Delta \W'| $ completes the proof.
 \end{proof}



\subsection{Simplifying Assumptions}\label{sec: simplify}
The next three lemmas show that to prove Theorem \ref{thm: main} it will be sufficient to assume that $f\geq 0$ and $f \in C^{\infty}$, and $f$ is compactly supported. The only quantitative assumption on $f$ will be that {$\| f \|_{L^{\infty}(\mathbb{R}^n)}\leq 1$.}

 \begin{lemma} \label{l:fnonneg}
   It is sufficient to prove the estimate \eqref{eqn: quant stability main} of  Theorem \ref{thm: main} in the case when  $f \geq 0$. 
 \end{lemma}

 \begin{proof}
 Choose any $f$ with $\| f\|_{L^{\infty}} \leq 1$ and  write $f = f^+ - f^-$. Note that $\| f_+\|_{L^{\infty}}, \| f_-\|_{L^{\infty}} \leq 1$
   \[
   u_{\W}^f = u_{\W}^{f^+} - u_{\W}^{f^-} \quad \text{ and } \quad 
   u_{B(x_{\W})}^f = u_{B(x_{\W})}^{f^+} - u_{B(x_{\W})}^{f^-}. 
   \]
   So, if the estimate \eqref{eqn: quant stability main}  is true whenever the right hand side is nonnegative with $L^{\infty}$ norm at most $1$, then applying it to $\W$ with $f_-$ and $f_+$,
   \[
   \begin{aligned}
   \int \big|u_{\W}^f- u_{B(x_{\W})}^f \big|^2
   &= \int \big|u_{\W}^{f^+} - u_{B(x_{\W})}^{f^+} -u_{\W}^{f^-} + u_{B(x_{\W})}^{f^-}  \big|^2\\
   &\leq 2\int \big|u_{\W}^{f^+} - u_{B(x_{\W})}^{f^+} \big|^2
   + 2\int \big|u_{\W}^{f^-} - u_{B(x_{\W})}^{f^-} \big|^2\leq 4 C_{n}\, (\lambda_1(\W) - \lambda_1(B)). 
   \end{aligned}
   \]
 \end{proof}

\begin{lemma} \label{l:fsmooth}
 It is sufficient to prove  the estimate \eqref{eqn: quant stability main} of  Theorem \ref{thm: main} in the case when  $f \in C^{\infty}(\R^n)$. 
\end{lemma}

 \begin{proof}
 Choose any $f$ with $\| f\|_{L^{\infty}} \leq 1$.
  Utilizing the standard mollification, if $f_{\epsilon} := f \ast \eta_{\epsilon}$, then 
  $\| f_{\epsilon} \|_{L^{\infty}(\mathbb{R}^n)} \leq \| f \|_{L^{\infty}(\mathbb{R}^n)} \leq 1$. Furthermore, $f_{\epsilon} \to f$ in $L^2(\mathbb{R}^n)$. Now 
  \begin{equation} \label{e:fepsilon}
  \begin{split}
    \int \big|u_{\W}^f- u_{B(x_{\W})}^f \big|^2
    &\leq \int \big|u_{\W}^f- u_{\W}^{f_{\epsilon}} \big|^2 + \int \big|u_{\W}^{f_{\epsilon}}- u_{B(x_{\W})}^{f_{\epsilon}} \big|^2  
    + \int \big|u_{B(x_{\W})}^{f_{\epsilon}}- u_{B(x_{\W})}^f \big|^2 
    \end{split}
  \end{equation}
To estimate the first term on the right-hand side of \eqref{e:fepsilon}, note that
  \[
   \begin{aligned}
       \| \nabla (u_{\W}^f- u_{\W}^{f_{\epsilon}}) \|^2_{L^2(\W)}
       &= \int_{\W} (f-f^{\e})(u_{\W}^f- u_{\W}^{f_{\epsilon}}) \\
       &\leq \| f - f^{\e} \|_{L^2(\W)} \| u_{\W}^f- u_{\W}^{f_{\epsilon}}\|_{L^2(\W)}. 
   \end{aligned}
  \]
On the other hand, from the Faber-Krahn inequality we have 
  \[
  \| u_{\W}^f - u_{\W}^{f_{\e}} \|^2_{L^2(\W)} \leq \frac{1}{\lambda_1(B)} \|\nabla (u_{\W}^f - u_{\W}^{f_{\e}}) \|^2_{L^2(\W)}.
  \]
Putting these together gives 
  \[
  \| u_{\W}^f - u_{\W}^{f_{\e}} \|_{L^2(\W)} \leq \frac{1}{\lambda_1(B)} \| f - f^{\e} \|_{L^2(\W)}. 
  \]
  The same estimate will also apply to the third term in \eqref{e:fepsilon}. 
  %
  So, if the estimate \eqref{eqn: quant stability main} of Theorem~\ref{thm: main} holds for smooth right-hand sides, then applying this estimate to $f^\e$ shows the second term in \eqref{e:fepsilon} is bounded above by $C_{n}(\lambda_1(\W) - \lambda_1(B))$.
  Letting $\e \to 0$ in \eqref{e:fepsilon}, we conclude the result.  
 \end{proof}



\begin{lemma}\label{lem:fcptspt}
    It is sufficient to prove  the estimate \eqref{eqn: quant stability main} of  Theorem \ref{thm: main} in the case when  $f$ has compact support.
\end{lemma}
%
%

 \begin{proof}
 Fix a smooth cutoff function $\psi:\R_+\to \R$ that is identically equal to $1$ in $(0,1)$ and vanishes on $(2,\infty)$ with $|\psi'|\leq 5$. 
Let $\W$ be an open bounded set with $|\W| = 1$.
  Choose any $f$ with $\| f\|_{L^{\infty}} \leq 1$.   
  
  Now, choose $R\geq 10$ such that $\W \subset B_{R/2}(0)$. Now, let 
 \[
{g(x)  := \psi\Big(\tfrac{|x|}{R}\Big) \,f(x)}
 \]
 By construction, $g$ is compactly supported and agrees with {$f$} in $B_R(0)$.
 
 Finally, by the choice of $R$ and definition of $x_\W$ (to be precise, by Proposition~\ref{lemma: qual stability} and \eqref{eqn: bary lip}), we see that $x_\W \subset B_{R/2}(0)$ and thus $B(x_{{\W}}) \subset B(0, R)$. 
 In particular, $g=f$ on $B({x_\W})\cup \W$, and thus by linearity of the Poisson equation \eqref{eqn: PDE with f RHS}, 
 \[
 \big\| u_\W^f - u_{B(x_\W)}^f \big\|_{L^2} =  \big\| u_\W^g - u_{B(x_\W)}^g \big\|_{L^2}\,.
 \]
 So, if the estimate \eqref{eqn: quant stability main} of Theorem~\ref{thm: main} holds for compactly supported right-hand sides, then the estimate applied with right-hand side $g$ immediately yields the estimate with right-hand side $f$. 
 \end{proof}



%
\section{Resolvent Estimates for Nearly Spherical Sets}\label{sec: linear stability}
 %
A set $\Omega\subset \R^n$ is called a {\it nearly spherical set} parametrized by $\phi \in C^{2,\gamma}(\partial B, \mathbb{R})$
 if $|\W|=\om_n$,  $x_\W =0$, and 
\[
\partial \Omega :=  \{ (1+\phi(x))x :  x \in \partial B\}\,
\]
for a function $\phi : \partial B\to \R$ with $\|\phi\|_{C^0}< 1.$ The truncated barycenter $x_\W$ of a nearly spherical set $\W$ agrees with the classical barycenter. We show Theorem~\ref{thm: main} holds if $\W$ is a nearly spherical set, provided $\|\phi\|_{C^{1,\alpha}}$ is small with $1/2< \alpha < 1$.
\begin{theorem}\label{cor: Fuglede}
	Fix {$\alpha \in (1/2,1)$}. There exist $\delta = \delta(n,\alpha)$  and $c=c(n,\alpha)$ such that if $f \in L^\infty(\R^n)$ with $\|f \|_{L^\infty}\leq 1$ and  $\W$ is a nearly spherical set of class $C^{2,\gamma}$ parametrized by $\phi$ with $\| \phi\|_{C^{1,\alpha}} \leq \delta$, then
\[
\tor(\W) - \tor(B) \geq c \int \big| u_\W^f - u_{B({x_\W})}^f\big|^2.  
\]
\end{theorem}  
A main tool in the proof is a spectral gap estimate.   

\begin{lemma} \label{thm: spectral gap}
Fix  $\alpha \in (1/2,1)$. There exists $\delta = \delta(n,\alpha)$ such that if $\W$ is a nearly spherical set of class $C^{2,\gamma}$ parametrized by $\phi$ with {$\| \phi\|_{C^{1,\alpha}} \leq \delta$}, 
\[
\tor(\W) - \tor(B) \geq \frac{1}{32n^2} \| \phi \|_{H^{1/2}(\partial B)}^2. 
\]
\end{lemma}

  {  If we assume that $\| \phi\|_{C^{2,\gamma}} \leq \delta$, then the above result is Theorem 3.3 in \cite{BDV15}. Since we only assume $f \in L^{\infty}$ and not necessarily H\"older continuous, we will need the weaker assumption that $\| \phi\|_{C^{1,\alpha}} \leq \delta$ for some $\alpha>1/2$. The proof with this weaker assumption is obtained using the ideas in the recent paper \cite{p24}, and we provide the details in the appendix. }



Theorem~\ref{cor: Fuglede} follows directly by combining 
Lemma~\ref{thm: spectral gap} and the following lemma.
\begin{lemma}\label{lem: beta on NSS}
{Fix $\alpha \in (1/2,1)$. There exist $\delta = \delta(n,\alpha)$ and $c=c(n,\alpha)$ such that if $f \in L^\infty(\R^n)$ with $\| f\|_{L^\infty}\leq 1$ and $\W$ is a nearly spherical set of class $C^{2,\gamma}$ parametrized by $\phi$ with $\| \phi\|_{C^{1,\alpha}} \leq \delta$,} then
	\[
	\| \phi \|_{H^{1/2}(\partial B)}^2 \geq c \int \big| u_\W^f - u_{B}^f\big|^2 .
	\]
\end{lemma}


Toward proving Lemma~\ref{lem: beta on NSS},  for $\phi: \partial B\to \R$ with {$\| \phi\|_{C^{1,\alpha}} \leq \delta$}, we consider the solution $h$ to 
\begin{equation} \label{e:h}
 \begin{cases}
  &\Delta h=0 \quad \text{ in } B \\
  &h=\phi \quad \text{ on } \partial B. 
 \end{cases}
\end{equation}
Recall that $\| h\|_{H^1(B_1)}=\| \phi\|_{H^{1/2}(\partial B_1)}$. 
From the maximum principle,  $\|h \|_{C^1(\overline{B})} \leq \| \phi \|_{C^1(\partial B)} 
\leq \| \phi\|_{C^{1,\alpha}(\partial B)}\leq \delta$. 



We define a diffeomorphism  $\Psi: B\to \Omega$  by  $\Psi(x) = (1 + h(x) ) x$. If $\delta>0$ is chosen to be small, then $D\Psi(x)$ is nowhere vanishing and $\Psi$ is indeed a diffeomorphism from $B$ onto its image. Moreover, $\Psi$ maps $\partial B$ to $\partial \W$ and thus the image of $B$ under $\Psi$ is $\W$. Under these assumptions, we have the following:
\begin{lemma} \label{l:LNbound}
 Let $f \in L^{\infty}(\mathbb{R}^n)$ with $\| f\|_{L^\infty}\leq 1$. Then  
 \[
 \| f- f\circ \Psi \|_{H^{-1}(B)} \leq 2n\| \phi\|_{H^{1/2}(\partial B)}\,.
 \]
\end{lemma}



\begin{proof}
 Let $u \in C_0^1(B_1)$. From Lemma \ref{l:fsmooth} we may assume $f \in C^1(\mathbb{R}^n)$. 
 For $0 \leq t \leq \epsilon$, define
 \[
 \mathcal{F}(t):= \int_{B} f(x+th(x)x) \,u(x) \ dx. 
 \]
 From the mean value theorem, there exists $t_0 \in (0,\epsilon)$ such that 
 \[
 \int_{B}[f(x+h(x)x)-f(x)]\, u(x) \ dx =\mathcal{F}(1)-\mathcal{F}(0)=\mathcal{F}'(t_0)= \int_{B}\sum_{i=1}^n f_{x_i}(x+t_0h(x)x)x_i h(x)u(x). 
 \]
 Since $u h \in C_0^1(B)$, we may use integration by parts to obtain 
 \[
 \int_{B}[f(x+h(x)x)-f(x)]u(x) \ dx =- \int_{B} \sum_{i=1}^n f(x+t_0h(x)x)(u_{x_i}(x)h(x)x_i+ u(x)h_{x_i}(x)x_i + h(x)u(x)). 
 \]
 Using now that $x_i,f \in L^{\infty}$ and H\"older's inequality we have 
 \[
 \left| \int_{B}[f(x+h(x)x)-f(x)]u(x) \ dx \right| \leq 2n \| f \|_{L^{\infty}} \| u \|_{H_0^1(B_1)} \| h\|_{H^1(B_1)} = 2n \| f \|_{L^\infty}  \| u \|_{H_0^1(B_1)}\| \phi\|_{H^{1/2}(\partial B_1)}. 
 \]
\end{proof}
Using Lemma~\ref{l:LNbound}, we can now prove Lemma~\ref{lem: beta on NSS}.
\begin{proof}[Proof of Lemma \ref{lem: beta on NSS}]	
	We define $\hat{u}^f_\W : B\to \R$ by $\hat{u}^f_\W (x) = u^f_\W(\Psi(x))$, and we will prove that 
	\begin{align}
	\label{est 1}	\int |u^f_{\W}(x) - \hat{u}^f_\W (x)|^2 \,dx \leq C \| \phi \|_{H^{1/2}(\partial B)}^2,\\
	\label{est 2}	\int | \hat{u}^f_\W(x) - u^f_B(x)|^2 \, dx \leq C \| \phi \|_{H^{1/2} (\partial B)}^2
	\end{align}
	from which the lemma follows. We note that since $\phi \in C^{1,\alpha}(\partial B)$ and $\| f\|_{L^\infty(\R^n)} \leq 1$, by \cite[Theorem 8.33]{GT}, we have 
    \begin{equation} \label{e:c1est}
       \| u^f_{\W} \|_{C^1(\overline{\W})} , \, \| \hat{u}^f_{\W} \|_{C^1(\overline{B})}
       \leq C(n,\alpha). 
    \end{equation}

 
	
	To prove \eqref{est 1}, we observe that for any $x \in B\cap \W$, we have
	\[
	|u^f_\W(x) - \hat{u}^f_\W(x)| =|u^f_\W(x) - u^f_\W(\Psi(x))|\leq \| u^f_\W\|_{C^1(\W)} | x- \Psi(x)| \le  \| u^f_\W\|_{C^1(\W)} |h(x)|.
	\]
	So, 
	\[
	\int_{\W \cap B}|u^f_\W(x) - \hat{u}^f_\W(x)|^2 \,dx \leq C\| u^f_\W\|^2_{C^1(\W)}  \int_{\W \cap B} h^2(x) \leq  C\| u^f_\W\|_{C^1(\W)}  \| \phi \|_{L^2(\partial B)}^2\,.
	\]
	 Next, for $x \in \W\setminus B$ we have $|u^f_\W(x) - \hat{u}^f_\W(x)| = u^f_\W(x) \leq C\| u^f_\W\|_{C^1(\W)} d(x, \partial \W)$, so that 
	\begin{align*}
		\int_{\W \setminus B} |u^f_\W(x) - \hat{u}^f_\W(x)|^2 \,dx  & \leq C\| u^f_\W\|^2_{C^1(\W)}  \int_{\W \setminus B} d(x, \partial \W)^2 \, dx\\
		 &= C\| u^f_\W\|^2_{C^1(\W)}  \int_{S^{n-1} \cap \{ \phi>0\}} \int_1^{1+ \phi} (1+\phi -r)^2 \,r^{n-1} dr \, d\theta \\
		&=C\| u^f_\W\|^2_{C^1(\W)}  \int_{S^{n-1} \cap \{ \phi>0\}} \int_{0}^\phi s^2 \, ds \, d\theta  \leq \e \| \phi\|_{L^2(\partial B)}^2.
	\end{align*}
	An analogous argument shows the same bound for the integral over $B\setminus \W$. Summing up, we arrive at \eqref{est 1}.
	
	
	Next we prove \eqref{est 2}. We begin by noting that the pulled back function $\hat{u}^f_\W$ solves the equation
	\[
	\begin{cases}
		-\mathcal{L} \hat{u}^f_\W = \text{det} (d\Psi(x))f(\Psi(x))  & \text{ in } B\\
		\hat{u}^f_\W =0 & \text{ on } \partial B, 
	\end{cases}
	\]
	where $\mathcal{L}u = \text{div}(A \nabla u)$ for $A = m\mathcal{A}$, $\mathcal{A}(x) = d\Phi(\Psi (x) ) (d\Phi (\Psi(x))^*$ and $m = 1/\sqrt{\det \mathcal A}=\det(d \Psi(x))$ and $\Phi = \Psi^{-1}$ with the standard estimates 
	\begin{equation} \label{e:estimates}
	|A -\text{Id}| \leq C(h +|\nabla h|) , \qquad |\det(d\Psi(x))-1| \leq C(h+ |\nabla h|). 
	\end{equation}
	Let $v = u^f_B - \hat{u}^f_\W$; note that $v=0$ on $\partial B$ and that $-\Delta v = (f-mf(\Psi(x))) - (\mathcal{L}-\Delta)\hat{u}^f_\W. $ By the Poincar\'{e} inequality on the ball, it suffices to estimate $\int |\nabla v|^2 \leq C \| \phi \|_{H^{1/2}(\partial B)}^2.$ We have
	\begin{align*}
		 \int_B |\nabla v|^2\,dx& = \int_B \nabla v \cdot ( \nabla u^f_B - A\nabla \hat{u}^f_\W ) + \int_B \nabla v \cdot (A- \text{Id} ) \nabla \hat{u}^f_\W \\ 
   &=\int_B v \left( f(x)-\det(d \Psi(x))f(\Psi(x))\right) + \int_B \nabla v \cdot (A- \text{Id} ) \nabla \hat{u}^f_\W = (I) + (II). 
	\end{align*}
 Now 
 \[
 (I) = \int_{B} v(x) (f(x)-f(\Psi(x))) \ dx + \int_{B}v(x)f(\Psi(x))(1-\text{det} (d \Psi(x)))\ dx= (Ia)+ (Ib). 
 \]
 From Lemma \ref{l:LNbound} we have 
 \[
 |(Ia)|\leq n \| f \|_{L^{\infty}} \| v \|_{H_0^1(B)} \| \phi \|_{H^{1/2}(\partial B)}
 \leq \epsilon \int_{B} |\nabla v|^2 + \frac{C}{\epsilon} \| f \|^2_{L^{\infty}} \| \phi \|^2_{H^{1/2}(\partial B)}. 
 \]
 Using the estimates in \eqref{e:estimates} we also have 
 \[
 |(Ib)| \leq C \int_{B} |v(x)| |f(\Psi(x))| (|h| + |\nabla h|) \ dx 
 \leq \epsilon \int_{B} v^2 + \frac{C}{\epsilon} \| f \|^2_{L^{\infty}} \| \phi \|^2_{H^{1/2}(\partial B)}. 
 \]
 Again using \eqref{e:estimates} as well as estimates on $\hat{u}^f_{\Omega}$ we obtain 
 \[
 |(II)| \leq C \| \hat{u}^f_{\Omega} \|_{C^1(B)} \int_{B} |\nabla v| (|h|+ |\nabla h|) \ dx 
 \leq \epsilon \int_{B} |\nabla v|^2 + \frac{C}{\epsilon}  \| \hat{u}^f_{\Omega} \|^2_{C^1(B)} \| \phi \|^2_{H^{1/2}(\partial B)}. 
 \]
 Using \eqref{e:c1est} this proves \eqref{est 2} and concludes the proof. 
 \end{proof}





  
 
\section{Proofs of the Main Theorems}\label{sec: proofs of main theorems}
This section is dedicated to showing how the two main theorems, Theorem~\ref{thm: main} and Theorem~\ref{thm: eigenfunction estimates}, follow from Theorem~\ref{thm: summary for minimizers} and Theorem~\ref{cor: Fuglede}. For the time being, we assume Theorem~\ref{thm: summary for minimizers} holds; this theorem will be established in sections~\href{https://arxiv.org/abs/2504.13053v2}{6} and \href{https://arxiv.org/abs/2504.13053v2}{7}. 

\subsection{Proof of Theorem~\ref{thm: main}}\label{ssec: proof of main}
We begin with Theorem~\ref{thm: main}, which, as explained in the introduction, follows from Theorem~\ref{thm: summary for minimizers} and Theorem~\ref{cor: Fuglede} using a classical selection principle argument.
\begin{proof}[Proof of Theorem~\ref{thm: main}]

 {\it Step 1: Reductions and setup.}
Together Lemmas~\ref{lem: KJ consequence}, \ref{l:fnonneg}, \ref{l:fsmooth}, and \ref{lem:fcptspt} show that to prove Theorem~\ref{thm: main}, it suffices to show that there is a positive constant $c_{n}>0$ depending on $n$ such that for any bounded open set $\W \subset \R^n$ with $|\W| = \omega_n$ and for any nonnegative function $f \in C^\infty_c(\R^n)$ with $\|f\|_{L^{\infty}(\R^n)}\leq 1$.
\begin{align}
\label{eqn: quant stability main 2}
\tor(\W ) - \tor(B)  & \geq c_{n}\int_{\R^n} \big|u_{\W}^f - u_{B(x_\W)}^f\big|^2 \,
\end{align}
Here $x_\W$ is the truncated  barycenter defined in \eqref{eqn: barycenter definition}. Since the integral on the right-hand side is bounded above by $2\om_n\|w_B\|_{L^\infty}^2 \leq \om_n/2n^2 \leq 1,$ the estimate \eqref{eqn: quant stability main 2} holds trivially when $\tor(\W ) - \tor(B)\geq \e$ for any $\e>0$ by choosing the constant $c_{n}$ to be small enough depending on $\e$. So, it suffices to show the theorem when  is  $\tor(\W ) - \tor(B)$ is small.

We use the short-hand notation $\beta_f(\W) = \| u_\W^f - u_{B(x_\W)}^f\|_{L^2(\R^n)}$. 	Suppose by way of contradiction that there are open bounded sets $\{\W_j\}$ in $\R^n$ with $|\W_j| = \omega_n$ and  smooth, nonnegative, compactly supported functions  $\{f_j\}$ with $\| f_j\|_{L^{\infty}}\leq 1$ such that $0< \tor(\W_j)-\tor(B_1)\to 0$ but 
	\begin{equation}\label{e: sp contra}
		\tor(\W_j ) - \tor(B_1) \leq \frac{1}{j}\, \beta_{f_j}^2(\W_j). 
	\end{equation}



 {\it Step 2: Replacement of $\W_j$.}
Fix $\delta = \delta(n, \alpha)$ according to Theorem~\ref{cor: Fuglede} and let $\tau \leq  \bar{\tau}_0(n,\alpha,\delta) = \bar{\tau}_0(n,\alpha)$ be chosen according to Theorem~\ref{thm: summary for minimizers}.  Let $a_j = \beta_{f_j}(\W_j)^2 \in (0, 1]$ and consider the functional
	\[
\mathcal{F}_j(U) =	\tor(U) + \tau \sqrt{a_j^2 + (a_j -\beta^2_{f_j}(U))^2}\,.
	\]
By Theorem~\ref{summary ex and reg theorem}, there is a minimizer  $U_j$ of the problem
    	\begin{equation}
		\inf \{ \mathcal{F}_j(U) \ : \ U \subset \R^n \text{ open bounded set}, \ |U| = \omega_n\},
	\end{equation} 
and $U_j -x_{U_j}$ is a nearly spherical set parametrized by a function $\phi_j \in C^{2,\alpha}( \partial B_1, \R)$ with $\|\phi_j\|_{C^{1,\alpha}}\leq \delta.$ So, by Theorem~\ref{thm: spectral gap} applied with $f = f_j(\cdot -x_{U_j})$, 
	\begin{equation}\label{eqn: sp fuglede}
		\tor(U_j ) - \tor(B_1) \geq c\, \beta_{f_j}^2(U_j). 
	\end{equation}

 {\it Step 3: Reaching a contradiction.}   
Taking $\W_j$ as a competitor for the minimality of $\mathcal{F}_j$ and applying \eqref{e: sp contra}, 
\begin{align*}
	\tor(U_j)  + \tau  \Big( \sqrt{a_j^2 + (a_j -\beta^2_{f_j}(U))^2} -a_j\Big)  \leq \tor(\W_j) \leq  \tor(B_1) + \frac{2}{j} a_j
\end{align*} 
The first inequality implies that $\tor(U_j) \leq \tor(\W_j)$, while the far left- and right-hand sides of the inequality and the Saint-Venant inequality guarantee that 
\[
\sqrt{1 + (1 -\beta^2_{f_j}(U_j)/a_j)^2} -1  \leq \frac{2}{j \,\tau } 
\]  
and thus $\beta_{f_j}(U_j)\geq 2 a_j= 2\beta_{f_j}(\W_j)^2$ for $j$ sufficiently large. Together these two observations and \eqref{e: sp contra} imply
	\begin{equation}
		\tor(U_j ) - \tor(B_1) \leq \frac{4}{j}\, \beta_{f_j}^2(U_j). 
	\end{equation}
	This contradicts \eqref{eqn: sp fuglede} when $j$ is large enough, completing the proof.
\end{proof}
%
%
%

\subsection{Proof of Theorem~\ref{thm: eigenfunction estimates}} Next we prove Theorem~\ref{thm: eigenfunction estimates} as an application of Theorem~\ref{thm: main}. Recall that by Lemma~\ref{lem: KJ consequence}, it suffices to replace the Faber-Krahn deficit $\lambda_1(\W)-\lambda_1(B)$ in Theorem~\ref{thm: eigenfunction estimates} with the Saint-Venant deficit $\tor(\W) -\tor(B)$, and that in section~\ref{ssec: proof of main} we proved that Theorem~\ref{thm: main} holds with $\tor(\W) -\tor(B)$ in place of $\lambda_1(\W)-\lambda_1(B)$.
Throughout this section we assume by translation that without loss of generality the truncated barycenter $x_\W$ of $\W$ is the origin.

Recall from the introduction that we denote $u_{\W,k}$ and $\lambda_{\W,k}$ as the $k$-th eigenfunctions and eigenvalues of $\Omega$, with the eigenfunctions extended globally to be defined on $\R^n$ and  normalized so that 
\begin{equation} \label{e:one}
 \| u_{\W,k} \|_{L^2(\W)} = 1 =  \| u_{B,k} \|_{L^2(\W)}. 
\end{equation}
If $k=1$, then we simply write $u_{\W}$ and $\lambda_{\W}$. 



We will utilize the following qualitative control on the eigenvalues whose proof is accomplished with standard geometric measure theory and concentration compactness techniques. 
\begin{proposition} \label{p:standard}
 Let $\W\subset\R^n$ be an open bounded set with $|\W|=\om_n$. For any $k \in \mathbb{N}$ and $\epsilon >0$, there exists $\delta_k >0$ such that if 
 \[
 \tor(\W)-\tor(B) < \delta_k, 
 \]
 then 
 \[
  |\lambda_k(\W) -\lambda_k(B)| \leq \epsilon.  
 \]
\end{proposition}
The recent result of Bucur-Lamboley-Nahon-Prunier described in the introduction is a sharp quantitative version of the qualitative Proposition~\ref{p:standard}.

To illustrate the method, we first prove Theorem~\ref{thm: eigenfunction estimates} for the first eigenfunctions. 

\begin{theorem} \label{t:first}
 Let $\W\subset \mathbb{R}^n$ be an open bounded set with $|\W|=|B|$ and $x_\W = 0$. There is a constant $c>0$ depending on dimension such that 
 \[
 \tor(\W)-\tor(B) \geq c \int_{\mathbb{R}^n} |u_{\W}-u_{B}|^2. 
 \]
\end{theorem}

\begin{proof}
We use the notation $\delta(\W)=\tor(\W)-\tor(B)$. 
As in the proof of Theorem~\ref{thm: main} above, it suffices to consider the case when $\delta(\W)$ is small.
From Theorem~\ref{thm: main} we have 
\[
 \| u_{\W}^f - u_{B}^f \|^2_{L^2(\mathbb{R}^n)} \leq C \delta(\W)
\]
 for any $f \in {L^\infty}(\R^n)$ with $\| f\|_{{L^\infty}(\R^n)} \leq 1.$
The eigenfunction $\lambda_B u_B$ is {bounded} (indeed Lipschitz continuous), and we will consider $f=\lambda_B u_{B}$, so the inequality above becomes 
\begin{equation} \label{e:quant}
 \| u_{\W}^f - u_{B}^f \|^2_{L^2(\mathbb{R}^n)} \leq C \|\lambda_B u_B \|_{L^{\infty}(\mathbb{R}^n)} \delta(\W)\,.
\end{equation}
We now write 
\[
u_{\W}^f = \sum_{k=1}^{\infty} a_k u_{\W,k}
\quad \text{ and } \quad  
u_{B}\big|_{\W} = \sum_{k=1}^{\infty} b_k u_{\W,k}. 
\]
Then from \eqref{e:quant} we have 
\begin{equation} \label{e:separate}
  \|u_B \|^2_{L^2(B \setminus \W)} +  \| u_{\W}^f - u_{B}\big|_{\W} \|^2_{L^2(\mathbb{R}^n)} =   \| u_{\W}^f - u_{B} \|^2_{L^2(\mathbb{R}^n)} \leq C \delta(\W). 
\end{equation}
 In particular, 
\begin{equation} \label{e:epsab}
    \sum_{k=1}^{\infty} (a_k - b_k)^2 \leq C \delta(\W). 
\end{equation}
Using that $-\Delta u_{\W}^f = \lambda_B u_B \chi_{\W}$ we obtain that for each $k \in \mathbb{N}$,
\[
\lambda_B b_k = \lambda_{\W,k} a_k.  
\]
Then from \eqref{e:epsab} we have 
\[
\sum_{k=1}^{\infty} b_k^2 \left(1- \frac{\lambda_B}{\lambda_{\W,k}} \right)^2 \leq C \delta(\W). 
\]
The spectral gap for the eigenvalues on the ball gives $\lambda_{B}/\lambda_{B,2} \leq c_n <1$. Combining with Proposition~\ref{p:standard}, if $\delta(\W)$ is small enough, then 
$\lambda_{B}/\lambda_{\Omega,2} \leq \tilde{c}_n <1$, and so 
\[
  \frac{\lambda_{B}}{\lambda_{\Omega,k}} \leq \tilde{c}_n <1 \quad \text{ for all } k \in \mathbb{N}. 
\]
Then for a new constant $C_1$ we have 
\[
\sum_{k=2}^{\infty} b_k^2 \leq C_1 \delta(\W). 
\]
Recalling that $\|u_B\|_{L^2(\mathbb{R}^n)}=1$, we have from \eqref{e:separate} that
\[
\sum_{k=1}^{\infty} b_k^2 = \| u_B \big|_\W \|^2_{L^2(\W)} \geq 1 - C\delta(\W). 
\]
Then 
\[
b_1^2 \geq 1 - (C+C_1) \delta(\W),
\]
and so 
\[
\begin{aligned}
\| u_{\W} - u_B \|_{L^2(\mathbb{R}^n)}^2 
&=\| u_{\W} - u_B \big|_{\W} \|_{L^2(\W)}^2 + \|u_B \|^2_{L^2(B \setminus \W)} \\
&=(1-b_1)^2 + \sum_{k=2}^{\infty} b_k^2 + \|u_B \|^2_{L^2(B \setminus \W)} \leq C_2 \delta(\W). 
\end{aligned}
\]
\end{proof}

\begin{thebibliography}{99}
\bibitem[1]{acampora2025sharpquantitativetalentisinequality} Paolo Acampora and Jimmy Lamboley, \emph{Sharp quantitative {T}alenti's inequality in particular cases}, arXiv 2503.07337 (2025).
\bibitem[2]{aac86} N.~Aguilera, H.~W. Alt, and L.~A. Caffarelli, \emph{An optimization problem with volume constraint}, SIAM J. Control Optim. \textbf{24} (1986), no.~2, 191--198. \MR{826512}
\bibitem[3]{AKN2} Mark Allen, Dennis Kriventsov, and Robin Neumayer, \emph{Linear stability implies nonlinear stability for {F}aber--{K}rahn-type inequalities}, Interfaces Free Bound. \textbf{25} (2023), no.~2, 217--324. \MR{4595296}
\bibitem[4]{AKN1} \bysame, \emph{Sharp quantitative {F}aber-{K}rahn inequalities and the {A}lt-{C}affarelli-{F}riedman monotonicity formula}, Ars Inven. Anal. (2023), Paper No. 1, 49. \MR{4570587}
\bibitem[5]{AKN3} \bysame, \emph{Rectifiability and uniqueness of blow-ups for points with positive {A}lt--{C}affarelli--{F}riedman limit}, Math. Ann. \textbf{391} (2025), no.~4, 6251--6289. \MR{4884572}
\bibitem[6]{aks23} Mark Allen, Dennis Kriventsov, and Henrik Shahgholian, \emph{The inhomogeneous boundary {H}arnack principle for fully nonlinear and {$p$}-{L}aplace equations}, Ann. Inst. H. Poincar\'e{} C Anal. Non Lin\'eaire \textbf{40} (2023), no.~1, 133--156. \MR{4565415}
\bibitem[8]{AC81} H.~W. Alt and L.~A. Caffarelli, \emph{Existence and regularity for a minimum problem with free boundary}, J. Reine Angew. Math. \textbf{325} (1981), 105--144. \MR{618549}
\bibitem[9]{amato2024talenticomparisonresultquantitative} Vincenzo Amato, Rosa Barbato, Alba~Lia Masiello, and Gloria Paoli, \emph{The talenti comparison result in a quantitative form}, 2024.
\bibitem[11]{ahmnt17} Jonas Azzam, Steve Hofmann, Jos\'e{}~Mar\'ia Martell, Kaj Nystr\"om, and Tatiana Toro, \emph{A new characterization of chord-arc domains}, J. Eur. Math. Soc. (JEMS) \textbf{19} (2017), no.~4, 967--981. \MR{3626548}
\bibitem[12]{Bhat01} Tilak Bhattacharya, \emph{Some observations on the first eigenvalue of the {$p$}-{L}aplacian and its connections with asymmetry}, Electron. J. Differential Equations (2001), No. 35, 15. \MR{1836803}
\bibitem[13]{BDF} V.~B\"{o}gelein, F.~Duzaar, and N.~Fusco, \emph{A quantitative isoperimetric inequality on the sphere}, Adv. Calc. Var. \textbf{10} (2017), no.~3, 223--265. \MR{3667048}
\bibitem[14]{BDS15} Verena B\"{o}gelein, Frank Duzaar, and Christoph Scheven, \emph{A sharp quantitative isoperimetric inequality in hyperbolic {$n$}-space}, Calc. Var. Partial Differential Equations \textbf{54} (2015), no.~4, 3967--4017. \MR{3426101}
\bibitem[15]{BrascoDePSurvey} Lorenzo Brasco and Guido De~Philippis, \emph{Spectral inequalities in quantitative form}, Shape optimization and spectral theory, De Gruyter Open, Warsaw, 2017, pp.~201--281. \MR{3681151}
\bibitem[16]{BDV15} Lorenzo Brasco, Guido De~Philippis, and Bozhidar Velichkov, \emph{Faber-{K}rahn inequalities in sharp quantitative form}, Duke Math. J. \textbf{164} (2015), no.~9, 1777--1831. \MR{3357184}
\bibitem[17]{BrascoPratelli} Lorenzo Brasco and Aldo Pratelli, \emph{Sharp stability of some spectral inequalities}, Geom. Funct. Anal. \textbf{22} (2012), no.~1, 107--135. \MR{2899684}
\bibitem[18]{BucurRobinLaplacian} Dorin Bucur, Alessandro Giacomini, and Paola Trebeschi, \emph{Stability results for the {R}obin-{L}aplacian on nonsmooth domains}, SIAM J. Math. Anal. \textbf{54} (2022), no.~4, 4591--4624. \MR{4461578}
\bibitem[19]{blnp23} Dorin Bucur, Jimmy Lamboley, Mickaël Nahon, and Raphaël Prunier, \emph{Sharp quantitative stability of the {D}irichlet spectrum near the ball}, arXiv 2304.10916 (2023).
\bibitem[20]{BuDalMaso} Giuseppe Buttazzo and Gianni Dal~Maso, \emph{An existence result for a class of shape optimization problems}, Arch. Rational Mech. Anal. \textbf{122} (1993), no.~2, 183--195. \MR{1217590}
\bibitem[21]{CaffSalsa} Luis Caffarelli and Sandro Salsa, \emph{A geometric approach to free boundary problems}, Graduate Studies in Mathematics, vol.~68, American Mathematical Society, Providence, RI, 2005. \MR{2145284}
\bibitem[22]{CiLe12} Marco Cicalese and Gian~Paolo Leonardi, \emph{A selection principle for the sharp quantitative isoperimetric inequality}, Arch. Ration. Mech. Anal. \textbf{206} (2012), no.~2, 617--643. \MR{2980529}
\bibitem[25]{DavidJerison} G.~David and D.~Jerison, \emph{Lipschitz approximation to hypersurfaces, harmonic measure, and singular integrals}, Indiana Univ. Math. J. \textbf{39} (1990), no.~3, 831--845. \MR{1078740}
\bibitem[27]{DET} Guy David, Max Engelstein, and Tatiana Toro, \emph{Free boundary regularity for almost-minimizers}, Adv. Math. \textbf{350} (2019), 1109--1192. \MR{3948692}
\bibitem[28]{DPMMCapacity} Guido De~Philippis, Michele Marini, and Ekaterina Mukoseeva, \emph{The sharp quantitative isocapacitary inequality}, Revista Matem\'{a}tica Iberoamericana (2021).
\bibitem[32]{Faber23} G.~Faber, \emph{Beweis, dass unter allen homogenen membranen von gleicher fl\"{a}che und gleicher spannung die kreisf\"{o}rmige den tiefsten grundton gibt}, Sitzungsber. Bayer. Akad. Wiss. M\"{u}nchen, Math.-Phys. Kl. (1923), 169--172.
\bibitem[33]{FiMaPr} A.~Figalli, F.~Maggi, and A.~Pratelli, \emph{A mass transportation approach to quantitative isoperimetric inequalities}, Invent. Math. \textbf{182} (2010), no.~1, 167--211. \MR{2672283}
\bibitem[34]{FTV2} I.~Fleschler, X.~Tolsa, and M.~Villa, \emph{Faber-{K}rahn inequalities, the {A}lt-{C}affarelli-{F}riedman formula, and {C}arleson's $\varepsilon^2$ conjecture in higher dimensions}, Preprint available at arXiv: 2306.06187 (2023).
\bibitem[35]{FTV1} Ian Fleschler, Xavier Tolsa, and Michele Villa, \emph{Carleson's $\varepsilon^2$ conjecture in higher dimensions}, 2023.
\bibitem[36]{FuJu} N.~Fusco and V.~Julin, \emph{A strong form of the quantitative isoperimetric inequality}, Calc. Var. Partial Differential Equations \textbf{50} (2014), no.~3-4, 925--937. \MR{3216839}
\bibitem[37]{FuMaPr} Nicola Fusco, Francesco Maggi, and Aldo Pratelli, \emph{The sharp quantitative isoperimetric inequality}, Ann. of Math. (2) \textbf{168} (2008), no.~3, 941--980. \MR{2456887}
\bibitem[38]{FMP09} \bysame, \emph{Stability estimates for certain {F}aber-{K}rahn, isocapacitary and {C}heeger inequalities}, Ann. Sc. Norm. Super. Pisa Cl. Sci. (5) \textbf{8} (2009), no.~1, 51--71. \MR{2512200}
\bibitem[39]{FuscoZhang} Nicola Fusco and Yi~Ru-Ya Zhang, \emph{A quantitative form of {F}aber-{K}rahn inequality}, Calc. Var. Partial Differential Equations \textbf{56} (2017), no.~5, Paper No. 138, 44. \MR{3691738}
\bibitem[40]{GT} David Gilbarg and Neil~S. Trudinger, \emph{Elliptic partial differential equations of second order}, Classics in Mathematics, Springer-Verlag, Berlin, 2001, Reprint of the 1998 edition. \MR{1814364}
\bibitem[41]{HansenNadir94} W.~Hansen and N.~Nadirashvili, \emph{Isoperimetric inequalities in potential theory}, Proceedings from the {I}nternational {C}onference on {P}otential {T}heory ({A}mersfoort, 1991), vol.~3, 1994, pp.~1--14. \MR{1266215}
\bibitem[42]{hp18} Antoine Henrot and Michel Pierre, \emph{Shape variation and optimization}, EMS Tracts in Mathematics, vol.~28, European Mathematical Society (EMS), Z\"urich, 2018, A geometrical analysis, English version of the French publication [MR2512810] with additions and updates. \MR{3791463}
\bibitem[43]{HMT} Steve Hofmann, Marius Mitrea, and Michael Taylor, \emph{Singular integrals and elliptic boundary problems on regular {S}emmes-{K}enig-{T}oro domains}, Int. Math. Res. Not. IMRN (2010), no.~14, 2567--2865. \MR{2669659}
\bibitem[44]{JTV21} Benjamin Jaye, Xavier Tolsa, and Michele Villa, \emph{A proof of {C}arleson's {$\varepsilon^2$}-conjecture}, Ann. of Math. (2) \textbf{194} (2021), no.~1, 97--161. \MR{4276285}
\bibitem[45]{JKNTA} David~S. Jerison and Carlos~E. Kenig, \emph{Boundary behavior of harmonic functions in nontangentially accessible domains}, Adv. in Math. \textbf{46} (1982), no.~1, 80--147. \MR{676988}
\bibitem[46]{MR4859580} Mikhail Karpukhin, Micka\"el Nahon, Iosif Polterovich, and Daniel Stern, \emph{Stability of isoperimetric inequalities for {L}aplace eigenvalues on surfaces}, J. Differential Geom. \textbf{129} (2025), no.~2, 415--490. \MR{4859580}
\bibitem[47]{k94} Carlos~E. Kenig, \emph{Harmonic analysis techniques for second order elliptic boundary value problems}, CBMS Regional Conference Series in Mathematics, vol.~83, Conference Board of the Mathematical Sciences, Washington, DC; by the American Mathematical Society, Providence, RI, 1994. \MR{1282720}
\bibitem[48]{KJ2} Marie-Th\'{e}r\`ese Kohler-Jobin, \emph{D\'{e}monstration de l'in\'{e}galit\'{e} isop\'{e}rim\'{e}trique {$P\lambda ^{2}\geq \pi j^{4}_{0}/2$}, conjectur\'{e}e par {P}\'{o}lya et {S}zeg\"{o}}, C. R. Acad. Sci. Paris S\'{e}r. A-B \textbf{281} (1975), no.~2-3, Aiii, A119--A121. \MR{385691}
\bibitem[49]{KJ1} \bysame, \emph{Une m\'{e}thode de comparaison isop\'{e}rim\'{e}trique de fonctionnelles de domaines de la physique math\'{e}matique. {I}. {U}ne d\'{e}monstration de la conjecture isop\'{e}rim\'{e}trique {$P\lambda ^{2}\geq \pi j^{4}_{0}/2$} de {P}\'{o}lya et {S}zeg\"{o}}, Z. Angew. Math. Phys. \textbf{29} (1978), no.~5, 757--766. \MR{511908}
\bibitem[50]{KJ3} \bysame, \emph{Symmetrization with equal {D}irichlet integrals}, SIAM J. Math. Anal. \textbf{13} (1982), no.~1, 153--161. \MR{641547}
\bibitem[51]{Krahn25} E.~Krahn, \emph{\"{U}ber eine von rayleigh formulierte minimaleigenschaft des kreises}, Math. Ann. \textbf{94} (1925), 97--100.
\bibitem[53]{MaggiBook} Francesco Maggi, \emph{Sets of finite perimeter and geometric variational problems}, Cambridge Studies in Advanced Mathematics, vol. 135, Cambridge University Press, Cambridge, 2012, An introduction to geometric measure theory. \MR{2976521}
\bibitem[54]{Melas} Antonios~D. Melas, \emph{The stability of some eigenvalue estimates}, J. Differential Geom. \textbf{36} (1992), no.~1, 19--33. \MR{1168980}
\bibitem[55]{Mukoseeva+2021} Ekaterina Mukoseeva, \emph{The sharp quantitative isocapacitary inequality (the case of p-capacity)}, Advances in Calculus of Variations (2021).
\bibitem[56]{N16} R.~Neumayer, \emph{A strong form of the quantitative {W}ulff inequality}, SIAM J. Math. Anal. \textbf{48} (2016), no.~3, 1727--1772. \MR{3498512}
\bibitem[57]{p24} Rapha\"el Prunier, \emph{Fuglede-type arguments for isoperimetric problems and applications to stability among convex shapes}, SIAM J. Math. Anal. \textbf{56} (2024), no.~2, 1560--1603. \MR{4711952}
\bibitem[59]{Talenti1976} Giorgio Talenti, \emph{Elliptic equations and rearrangements}, Annali della Scuola Normale Superiore di Pisa - Classe di Scienze \textbf{3} (1976), no.~4, 697--718 (eng).
\end{thebibliography}
\end{document}
